% source-sha256: 7db9d60041e090b7f53ee5d05947ec50ee8ce80561c5aeeace2f354bfcd674f5
% Options for packages loaded elsewhere
\PassOptionsToPackage{unicode}{hyperref}
\PassOptionsToPackage{hyphens}{url}
\documentclass[
]{article}
\usepackage{xcolor}
\usepackage[letterpaper,margin=1in]{geometry}
\usepackage{amsmath,amssymb}
\setcounter{secnumdepth}{-\maxdimen} % remove section numbering
\usepackage{iftex}
\ifPDFTeX
  \usepackage[T1]{fontenc}
  \usepackage[utf8]{inputenc}
  \usepackage{textcomp} % provide euro and other symbols
\else % if luatex or xetex
  \usepackage{unicode-math} % this also loads fontspec
  \defaultfontfeatures{Scale=MatchLowercase}
  \defaultfontfeatures[\rmfamily]{Ligatures=TeX,Scale=1}
\fi
\usepackage{lmodern}
\ifPDFTeX\else
  % xetex/luatex font selection
\fi
% Use upquote if available, for straight quotes in verbatim environments
\IfFileExists{upquote.sty}{\usepackage{upquote}}{}
\IfFileExists{microtype.sty}{% use microtype if available
  \usepackage[]{microtype}
  \UseMicrotypeSet[protrusion]{basicmath} % disable protrusion for tt fonts
}{}
\makeatletter
\@ifundefined{KOMAClassName}{% if non-KOMA class
  \IfFileExists{parskip.sty}{%
    \usepackage{parskip}
  }{% else
    \setlength{\parindent}{0pt}
    \setlength{\parskip}{6pt plus 2pt minus 1pt}}
}{% if KOMA class
  \KOMAoptions{parskip=half}}
\makeatother
\usepackage{longtable,booktabs,array}
\usepackage{caption}
\captionsetup[table]{skip=6pt}
\newcounter{none} % for unnumbered tables
\usepackage{calc} % for calculating minipage widths
% Correct order of tables after \paragraph or \subparagraph
\usepackage{etoolbox}
\makeatletter
\patchcmd\longtable{\par}{\if@noskipsec\mbox{}\fi\par}{}{}
\makeatother
% Allow footnotes in longtable head/foot
\IfFileExists{footnotehyper.sty}{\usepackage{footnotehyper}}{\usepackage{footnote}}
\makesavenoteenv{longtable}
\setlength{\emergencystretch}{3em} % prevent overfull lines
\providecommand{\tightlist}{%
  \setlength{\itemsep}{0pt}\setlength{\parskip}{0pt}}
% Release-document layout safeguards injected by scripts/generate_docs.sh.
\usepackage{fvextra}
\fvset{breaklines=true,breakanywhere=true,fontsize=\small}
\RecustomVerbatimEnvironment{verbatim}{Verbatim}{breaklines=true,breakanywhere=true,fontsize=\small}
\setlength{\emergencystretch}{3em}
\AtBeginDocument{\sloppy}
% Workflow-figure support: mermaid blocks in the canonical Markdown are mapped
% to hand-authored TikZ figures by scripts/pandoc_bounded_tables.lua.
\usepackage{tikz}
\usetikzlibrary{shapes.geometric,arrows.meta,positioning,fit,backgrounds,calc}
\definecolor{bdiStep}{HTML}{EAF2FB}
\definecolor{bdiEdge}{HTML}{365F91}
\definecolor{bdiText}{HTML}{17253A}
\definecolor{bdiGate}{HTML}{FFF5D6}
\definecolor{bdiGateEdge}{HTML}{9B6B00}
\definecolor{bdiOpt}{HTML}{F0EAF7}
\definecolor{bdiOptEdge}{HTML}{5B3E86}
\definecolor{bdiOut}{HTML}{E6F4EA}
\definecolor{bdiOutEdge}{HTML}{1E6B3A}
\tikzset{
  bdi/.style={draw=bdiEdge,line width=1pt,fill=bdiStep,text=bdiText,rounded corners=3pt,
              align=center,inner sep=4pt,font=\scriptsize\sffamily},
  bdigate/.style={bdi,draw=bdiGateEdge,fill=bdiGate},
  bdiopt/.style={bdi,draw=bdiOptEdge,fill=bdiOpt,dashed},
  bdiout/.style={bdi,draw=bdiOutEdge,fill=bdiOut},
  bdiflow/.style={-{Stealth[length=2mm]},draw=bdiEdge,line width=0.9pt},
  bdiback/.style={bdiflow,dashed},
  bdilbl/.style={font=\tiny\sffamily,text=bdiText,inner sep=1.5pt,fill=white},
}
\usepackage{bookmark}
\IfFileExists{xurl.sty}{\usepackage{xurl}}{} % add URL line breaks if available
\urlstyle{same}
% fallback for those not using the hyperref driver hyperxmp:
\makeatletter
\@ifundefined{xmpquote}{\newcommand{\xmpquote}[1]{#1}}{}
\makeatother
\hypersetup{
  pdftitle={Business Document Ingestion{,} Extraction \& Analytics Program},
  pdfauthor={Christopher Melhauser and theonlymuffinbot},
  hidelinks,
  pdfcreator={LaTeX via pandoc}}

\title{Business Document Ingestion, Extraction \& Analytics Program}
\usepackage{etoolbox}
\makeatletter
\providecommand{\subtitle}[1]{% add subtitle to \maketitle
  \apptocmd{\@title}{\par {\large #1 \par}}{}{}
}
\makeatother
\subtitle{Self-Contained Operating Plan and Vendor-Neutral Handoff
Specification --- Release 1.0}
\author{Christopher Melhauser and theonlymuffinbot}
\date{September 3, 2026}

\begin{document}
\maketitle

\section{0. Document Purpose and
Scope}\label{document-purpose-and-scope}

\subsection{0.0.1 Authorship and AI collaboration
context}\label{authorship-and-ai-collaboration-context}

Courtesy credit for this manuscript is given to Christopher Melhauser
and \texttt{theonlymuffinbot}. In this context,
\texttt{theonlymuffinbot} describes collaborative AI tooling. This
describes the tooling collaboration and does not assign legal authorship
or ownership to an AI system; legal attribution remains with human
contributors.

\subsection{0.1 What this document is}\label{what-this-document-is}

This is a self-contained operating plan for converting retained business
documents --- invoices, receipts, purchase orders, memos, remittances,
statements, and shipping records where present --- into a validated,
reconciled financial and operational dataset suitable for an agreed
business-system handoff.

ACKs, jobs, commission statements, shipment records, black-and-white
scans, sales-origin analysis, and email are illustrative configuration
examples only. They are not assumed client facts, required fields, or
delivery commitments. Replace them with the engagement's approved
document families, relationship keys, controls, and business questions
during Phase 0.

\subsection{0.2 Delivery package}\label{delivery-package}

This plan is self-contained for client delivery: it defines the
operating gates, evidence-preservation rules, quality controls, data
model, output contract, and remaining integration limits. The current
delivery package preserves source PDFs and run artifacts and supplies
the final review package as canonical JSON plus CSV, formatted XLSX
workbook, and HTML reviewer views.
\texttt{references/\allowbreak{}artifact-\allowbreak{}contracts.md}
supplies the definitive review interface;
\texttt{references/\allowbreak{}extraction-\allowbreak{}schema.md} and
\texttt{references/\allowbreak{}derived-\allowbreak{}field-\allowbreak{}schema.md}
supply the field dictionaries. The \texttt{examples/\allowbreak{}}
folder is explicitly fictional field-shape material, not an operational
review package.

\subsection{0.2.1 Delivery workflow and optional AI
integration}\label{delivery-workflow-and-optional-ai-integration}

\textbf{Source PDFs, references, and operating constraints} →
\textbf{preserve/profile} → \textbf{immutable page intake} →
\textbf{independent proposals and consensus} → \textbf{arithmetic,
validation, attribution, completeness, and sampling} →
\textbf{exhaustive final client review} → \textbf{analytics and agreed
load handoff}.

Optional OpenAI, Gemini-on-Vertex, OpenRouter, Document AI, HTR,
semantic-discovery, address, and LLM-adjudication lanes feed retained
evidence or proposals into this same path. None may bypass the
deterministic controls, close an exception, or supply client approval.

The selected run-level LLM adapter does not control the final decision
path. It runs only after immutable intake, retains raw provider evidence
and a page hash, and contributes one structured proposal to consensus.
Provider uncertainty, type proposals, handwriting flags, and failures
enter the review gate; no adapter can approve a record or close an
exception. Primary and secondary consensus lanes are invoked separately
and count as independent only when their providers are genuinely
independent.

The separately disabled LLM adjudication lane is intentionally narrower.
It may query only explicit nonfinancial, printed candidates after clear
deterministic validation and exact independent-extractor agreement.
\texttt{LLM\_\allowbreak{}ADJUDICATION\_\allowbreak{}MIN\_\allowbreak{}CONFIDENCE}
defaults to \texttt{0.99} and can be adjusted in \texttt{.env} only
through \texttt{1.0}; it cannot override financial, address,
handwriting, reassembly, or disagreement controls. Every result is an
audit-marked amendment proposal or exception and always enters the final
client-review gate.

\subsection{0.2.2 Semantic schema discovery, entity history, and future
CRM/RAG
delivery}\label{semantic-schema-discovery-entity-history-and-future-crmrag-delivery}

New source layouts are governed by a versioned canonical vocabulary and
a source-alias registry.
\texttt{scripts/\allowbreak{}schema\_\allowbreak{}discovery.py} produces
strict LLM suggestions for canonical field, semantic type, alternatives,
confidence, rationale, and supplied evidence only. It cannot create a
canonical field, alter a party or entity master, or approve a mapping.
Each new alias/template requires explicit client approval; the registry
update creates a new effective-dated snapshot. Familiar sources
presenting changed headers, dealer or brand names, city clues, or
relationship evidence return to review. Names, addresses, and
relationships are temporal evidence, not overwritten labels.

Use address evidence from the source page and the separate
address-validation stage first. Google Places (New), when enabled, is
only a bounded candidate-discovery source for observed dealer/brand
names and city hints. It never establishes entity identity, selects a
city/address, or bypasses review. The approved registry and validated
provenance-linked records constitute a vendor-neutral CRM staging
interface. The implemented local retrieval database indexes approved
canonical facts and preserves document/page/source-hash and
mapping-version citations. A specific CRM connector, remote vector
service, or ChatGPT connection remains deferred until the client selects
its security, retention, and access model.

\subsection{0.2.3 Flexible allocation
policy}\label{flexible-allocation-policy}

Commission and sales-credit logic is a versioned policy layer, not a
globally hard-coded percentage or a column rename.
\texttt{allocation\_\allowbreak{}policy.py} calculates effective
commission rate from retained source amounts and separately records
stated rates, shares, and formula proof. A client-approved rule may use
any approved source dimension---brand, dealer, product, client,
location, or a later schema-approved layer---and an effective-date
window. The most-specific active rule applies; missing/ambiguous rules,
unknown layers, invalid dates, and formula mismatches remain review
work.

For a new report layout, strict LLM output proposes allocation roles and
a policy card only. The client edits only the generated decision return
file to approve/defer, specify dimensions and dates, and make an
explicit approved override. The next registry snapshot preserves the
prior policy history. The LLM cannot create sales credit or alter a
historical result. See
\texttt{references/\allowbreak{}allocation-\allowbreak{}policy.md}.

The implemented control layer is ready to operate when client records,
reference exports, and approved provider configurations are received. It
does not claim that OCR/HTR vendor calls, handwriting-region detection,
broad automatic classification, or CRM loading have been deployed; those
require client-approved integrations and a representative golden-set
evaluation.

\subsection{0.3 Capability status and future
steps}\label{capability-status-and-future-steps}

This plan retains the 1.0.0 control baseline and also describes
implemented development additions in this checkout. See
\texttt{HANDOFF.md} for merge/deployment state; an \textbf{Unreleased}
capability is not a claim of production acceptance.

{\def\LTcaptype{none} % do not increment counter
\begin{longtable}[]{@{}
  >{\raggedright\arraybackslash}p{(\linewidth - 4\tabcolsep) * \real{0.3333}}
  >{\raggedright\arraybackslash}p{(\linewidth - 4\tabcolsep) * \real{0.3333}}
  >{\raggedright\arraybackslash}p{(\linewidth - 4\tabcolsep) * \real{0.3333}}@{}}
\toprule\noalign{}
\begin{minipage}[b]{\linewidth}\raggedright
Capability
\end{minipage} & \begin{minipage}[b]{\linewidth}\raggedright
Status
\end{minipage} & \begin{minipage}[b]{\linewidth}\raggedright
Engagement next step
\end{minipage} \\
\midrule\noalign{}
\endhead
\bottomrule\noalign{}
\endlastfoot
Immutable intake, proposal-only extraction, consensus, arithmetic,
validation, review packages, and JSON graph overlays & Implemented
control layer & Select the applicable source records and run
configuration. \\
Configured-primary/configured-buddy relationship inference and
cross-packet verification & Implemented proposal workflow & Set an
approved provider budget; review retained proposals and exceptions. \\
Reference schema, attribution keys, analytic dimensions, and allocation
policy & Configurable implemented controls & Approve the relationship
key, field mappings, and policy for the engagement. \\
Approved-fact MCP/API, account cards, governed queries, sales slices,
seven standard reports and remote CSV/XLSX downloads & Implemented
shared service and transports & Build an approved non-empty snapshot;
accept the selected TLS/OAuth/client deployment. \\
Schema-guided PNG/JPEG intake and local source-only export & Implemented
optional proposal pilot & Authorize image/model data flow; verify real
transfer, receipts, source package and normal pipeline re-entry. \\
Authoritative GL, payment, reference, or authorization evidence &
External operating control & Supply it when available; otherwise retain
an unavailable-control exception. \\
Additional OCR/HTR vendors, calibrated detector/template deployment,
email ingestion, target CRM connector, hosted retrieval, and production
accuracy claims & Future integration work & Define vendor, security,
retention, evaluation, and acceptance criteria before implementation. \\
\end{longtable}
}

Nothing in this template turns an inference or a missing client control
into an approved fact. Future steps remain explicitly future until their
integration, test evidence, and client acceptance are recorded.

\subsubsection{Optional visual intake and remaining adapter
work}\label{optional-visual-intake-and-remaining-adapter-work}

The current MCP servers can add eleven operations for schema discovery,
session discovery, original-page upload/read, proposal submission,
reviewer grants, review summaries, and status/proposal retrieval. The
remote OAuth service exposes equivalent JSON API routes. Intake is
disabled unless \texttt{-\/-ingestion-dir} is supplied. Local discovery
grows from 12 to 23 tools; remote discovery grows from 14 to at most 25,
filtered by scopes. The TLS/bearer GET API and default local launcher
remain unchanged. The combined servers still require an approved
snapshot and separate activation authorization.

The connected vision-capable assistant reads retained images; the server
makes no additional model call. Originals, source references, unknown
fields, uncertainties and all rejected/corrected proposals remain
append-only in a separate journal. A valid submission is pending review,
not an approved record. No native capture client or chat-attachment
relay is included. Client upload and vision support must be demonstrated
rather than inferred from MCP access.

\texttt{visual\_\allowbreak{}ingestion\_\allowbreak{}export.py}
snapshots an existing session read-only into a new source package,
preserves every original attempt/proposal/rejection, and only for
complete usable pages prepares a session-named image PDF with a
provenance map. Verify against the separately retained receipt, inspect
the PDF, and use normal profiling/intake. Missing or failed pages block
progression. Every source-package exception belongs in final review. The
package is neither a consensus handoff nor an approved CRM import.

The \href{../references/visual-ingestion.md}{visual contract} and
\href{TECHNICAL_DOCUMENTATION.md}{technical manual} define flags, sizes
and recovery. Governed create/update, shared review, saved custom
reports, safe general joins and broader analytics remain in the
\href{../references/business-data-platform-roadmap.md}{vendor-neutral
roadmap}. They do not depend on Salesforce. The existing no-send
common-object package already prepares twelve CRM import files; live
target adapters still require tenant-specific mapping, sandbox evidence
and separate authorization.

\subsection{0.4 What is in scope}\label{what-is-in-scope}

Ingestion and normalization; document reassembly, classification, and
field extraction; handwriting capture and linkage; \textbf{attribution
of every monetary amount to a configured business reference}; origin
resolution; validation, entity resolution, and reconciliation;
completeness controls covering all purchases and payments; analytics;
canonical data model and CRM load sequence. Email integration is in
scope as a deferred option.

\subsection{0.5 What is out of scope}\label{what-is-out-of-scope}

Construction of the CRM application itself. This program terminates at a
documented, validated, load-ready dataset plus the load order and
integrity rules the CRM must honour.

\subsection{0.6 Stated outcome}\label{stated-outcome}

\begin{enumerate}
\def\labelenumi{\arabic{enumi}.}
\tightlist
\item
  A \textbf{complete, auditable financial record} of every purchase and
  payment evidenced in the corpus, each figure traceable to a specific
  source page.
\item
  \textbf{Every monetary amount attributed} to a configured business
  reference, or explicitly registered as unattributable with a stated
  reason and a dollar value.
\item
  A \textbf{quantified statement of accuracy and completeness} --- a
  measured error rate with confidence bounds and an explicit enumeration
  of what is missing.
\item
  A \textbf{canonical dataset}, deduplicated and entity-resolved, ready
  for CRM load in a defined order.
\item
  A \textbf{standing analytics layer} answering the client's decision
  questions, including revenue by originating city and by job.
\item
  A \textbf{reusable pipeline} and a permanent regression test set.
\end{enumerate}

\subsection{0.7 Governing principles}\label{governing-principles}

\textbf{Nothing is silently discarded.} Every input page produces an
output record, even if that record is ``unreadable, escalated.''

\textbf{Original values are never overwritten.} Corrections are stored
as amendments alongside originals, with source attribution and effective
time.

\textbf{Documents prove themselves arithmetically.} A confidence score
is not evidence. Internal arithmetic consistency is. This principle
carries disproportionate weight in every financial-document engagement
--- see §1.2.

\textbf{Completeness is reported on the face of every analysis.} Never
plug a variance to make totals agree.

\textbf{Every dollar carries an attribution or an explicit reason it has
none.} Silently unattributed revenue defeats the purpose of the
exercise.

\begin{center}\rule{0.5\linewidth}{0.5pt}\end{center}

\section{1. Engagement Configuration and Open
Variables}\label{engagement-configuration-and-open-variables}

\subsection{1.1 Example: black-and-white scan
corpus}\label{example-black-and-white-scan-corpus}

The following is an illustrative configuration for a black-and-white
corpus. For each client, Phase 0.5 measures the actual colour depth, bit
depth, and scan quality; it does not assume this example or choose a
processing path before that evidence exists.

\subsubsection{Direct consequences}\label{direct-consequences}

{\def\LTcaptype{none} % do not increment counter
\begin{longtable}[]{@{}
  >{\raggedright\arraybackslash}p{(\linewidth - 4\tabcolsep) * \real{0.3333}}
  >{\raggedright\arraybackslash}p{(\linewidth - 4\tabcolsep) * \real{0.3333}}
  >{\raggedright\arraybackslash}p{(\linewidth - 4\tabcolsep) * \real{0.3333}}@{}}
\toprule\noalign{}
\begin{minipage}[b]{\linewidth}\raggedright
Capability
\end{minipage} & \begin{minipage}[b]{\linewidth}\raggedright
Status
\end{minipage} & \begin{minipage}[b]{\linewidth}\raggedright
Consequence
\end{minipage} \\
\midrule\noalign{}
\endhead
\bottomrule\noalign{}
\endlastfoot
Ink/colour separation for handwriting detection & \textbf{Unavailable} &
Pen and print are the same colour. Annotation detection falls back to
stroke shape and layout deviation. \\
Stamp and seal detection & \textbf{Unavailable by colour} & ``PAID''
stamps, approval stamps, and received-date stamps no longer separate
from print. Planned template matching requires client calibration
(§3.5). \\
Highlighter and markup detection & \textbf{Lost} & Highlighted regions
render as grey wash or, on bilevel scans, erase the underlying text
entirely. Affected regions are flagged as damaged, not silently read. \\
Handwriting recognition accuracy & \textbf{Degraded} & Lower stroke
fidelity, particularly on 1-bit scans. \\
Printed-text OCR & \textbf{Largely unaffected} & Print OCR is designed
for bilevel input. \\
Classification, matching, reconciliation, attribution &
\textbf{Unaffected} & Structural, not tonal. \\
\end{longtable}
}

\subsubsection{What this changes in the
plan}\label{what-this-changes-in-the-plan}

For a black-and-white corpus, use these compensating controls:

\begin{enumerate}
\def\labelenumi{\arabic{enumi}.}
\tightlist
\item
  \textbf{Use at least two genuinely independent extraction providers;
  target three where the client risk assessment requires triple entry.}
  Release 1.0 supplies the consensus contract and optional OpenAI,
  Gemini-on-Vertex, OpenRouter, and Document AI evidence adapters, but
  the client must select and accept the independent production
  providers. Single-engine error rates rise materially on bilevel input.
\item
  \textbf{Multi-binarization ensemble} on any 8-bit grayscale pages
  (§3.4). Each variant is retained comparison evidence; readings from
  one vendor remain correlated and are not separate independent engines.
\item
  \textbf{Arithmetic self-proof is promoted from a validation rule to
  the primary error-detection mechanism} (§1.2).
\item
  \textbf{Handwriting is assumed present until proven absent.}
  Proposal/detector output sets processing priority, never eligibility;
  local detection still requires client calibration.
\item
  \textbf{QA oversampling on handwriting-bearing strata rises to
  8--10×}.
\end{enumerate}

\subsubsection{The one remaining question about the
scans}\label{the-one-remaining-question-about-the-scans}

The corpus is black-and-white, but \textbf{8-bit grayscale and 1-bit
bilevel are very different inputs} and the difference is not visible to
the eye at normal zoom. Grayscale retains stroke intensity information
that stroke-morphology detection depends on; bilevel has discarded it
irreversibly.

Phase 0.5 measures the split. If a material share is bilevel, and
grayscale or colour masters still exist anywhere --- original paper, an
earlier scan generation, a scanner's retained output folder ---
re-deriving those pages from the better source is by far the cheapest
accuracy improvement available in this project.

\subsection{1.2 Why arithmetic proof now carries the
program}\label{why-arithmetic-proof-now-carries-the-program}

With ink separation gone, some handwritten annotations \textbf{will} go
undetected. That is a certainty, not a risk to be mitigated to zero.

The recovery mechanism is arithmetic. A handwritten quantity correction
or price override that the detector missed will, in most cases, break
the document's internal arithmetic --- the printed lines will no longer
sum to the printed total once the real quantity is what governs. The
failure surfaces the annotation that detection missed.

This is why the self-proof rule remains non-negotiable whenever
handwritten financial changes may be present. \textbf{Arithmetic proof
is the backstop when recognition is unavailable, deferred, or
intentionally out of scope.}

\subsection{1.3 Example configuration
parameters}\label{example-configuration-parameters}

{\def\LTcaptype{none} % do not increment counter
\begin{longtable}[]{@{}
  >{\raggedright\arraybackslash}p{(\linewidth - 2\tabcolsep) * \real{0.5000}}
  >{\raggedright\arraybackslash}p{(\linewidth - 2\tabcolsep) * \real{0.5000}}@{}}
\toprule\noalign{}
\begin{minipage}[b]{\linewidth}\raggedright
Parameter
\end{minipage} & \begin{minipage}[b]{\linewidth}\raggedright
Value
\end{minipage} \\
\midrule\noalign{}
\endhead
\bottomrule\noalign{}
\endlastfoot
Scan colour depth & Example: black-and-white source corpus; Phase 0.5
records the measured grayscale/bilevel split per page \\
QA model & Sampling-only on the clean population; 100\% review of all
exceptions \\
Efficiency posture & Robustness prioritised over cost \\
Financial coverage & All purchases and payments documented and
reconciled \\
Handwriting & Example: present; recognition deferred; capture-and-link
retained (§7) \\
Dollar attribution & Example: every dollar ties to an approved
relationship key such as ACK, job, project, order, or PO \\
Sales origin & Example analytic dimension; replace with approved
reporting dimensions \\
Email corpus & Optional future phase if separately authorized \\
Corpus period & Engagement-defined \\
Corpus ordering & Unordered; pages not grouped into documents \\
\end{longtable}
}

\subsection{1.4 Parameters requiring client
input}\label{parameters-requiring-client-input}

\begin{itemize}
\tightlist
\item
  Total page count across all scan files
\item
  \textbf{ACK and job-number format conventions}, including any changes
  over the five-year period
\item
  \textbf{Definition of ``sales origin city''} --- see §8.4, which
  explains why this must be settled explicitly
\item
  Materiality threshold for monetary-unit sampling
\item
  Tolerable error rate per field class
\item
  GL variance tolerance
\item
  Availability of grayscale or colour masters for any bilevel pages
\item
  Email export format, date range, and mailbox scope (if Phase 7
  proceeds)
\item
  Data residency and PII handling constraints
\end{itemize}

\begin{center}\rule{0.5\linewidth}{0.5pt}\end{center}

\section{2. Phase 0 --- Discovery}\label{phase-0-discovery}

Working session with the client. Output is a one-page constraint
document.

\subsection{2.1 Inventory}\label{inventory}

File count, total page count, scan DPI, single- or double-sided capture,
presence of any native text layer.

\subsection{2.2 Document taxonomy}\label{document-taxonomy}

Confirm which types exist. Do not assume the list. Candidates:
commercial invoice, packing list, bill of lading, air waybill, carrier
freight bill, customs entry, proof of delivery, credit memo, debit memo,
remittance advice, purchase order, \textbf{acknowledgement (ACK)},
statement of account.

The acknowledgement document type is load-bearing when it is a
configured relationship key. If ACKs exist as their own document class
in the corpus, they are the authoritative source for the attribution key
and are processed first.

\subsection{2.3 Attribution conventions}\label{attribution-conventions}

An essential discovery item when acknowledgement or order references are
in scope.

\begin{itemize}
\tightlist
\item
  What does an ACK number look like? Fixed prefix, length, check digit?
\item
  What does a job or project number look like?
\item
  Did either convention change during the measured period? Format
  changes are the most common cause of silent attribution failure.
\item
  Which is authoritative when a document carries both?
\item
  Are there dollars that legitimately have neither --- internal
  transfers, rebates, rebilled freight, credits? These need a defined
  disposition rather than showing up as failures.
\item
  Where does the number physically appear on each document type? Printed
  field, stamped, or handwritten?
\end{itemize}

The last question determines how much of §7 is required. If ACK numbers
are routinely handwritten, handwriting capture is not optional
regardless of the general descoping instruction.

\subsection{2.4 Systems of record}\label{systems-of-record}

Accounting platform, TMS, order-entry or quoting system, carrier
portals.

\textbf{An order-entry or quoting system is material}: when an
authoritative digital list of relationship keys, dates, parties, and
amounts exists, it becomes the attribution reference table and makes §8
more reliable than deriving everything from documents alone.

\subsection{2.5 Retention, privacy,
residency}\label{retention-privacy-residency}

What may leave the client network, permitted cloud regions, statutory
retention on originals. Extend to mailbox scope if Phase 7 is
anticipated.

\subsection{2.6 Decision list}\label{decision-list}

The questions the analytics must answer, as literal sentences. Defines
``finished.'' Include the city-of-sale question explicitly, because its
phrasing determines the §8.4 resolution hierarchy.

\subsection{2.7 Deliverables}\label{deliverables}

Constraint document; 50-page random calibration sample; ACK/job format
specification.

\begin{center}\rule{0.5\linewidth}{0.5pt}\end{center}

\section{3. Phase 0.5 --- Scan
Profiling}\label{phase-0.5-scan-profiling}

Runs against the \textbf{full corpus}, not the sample. About an hour of
compute.

The purpose is to measure the source evidence before selecting a
processing path. The profiling now characterises \textbf{how degraded
the black-and-white corpus is} and routes pages to the appropriate
enhancement path.

\subsection{3.1 Probe}\label{probe}

Per page: bit depth, effective DPI, compression codec, ink coverage
ratio, estimated blur, contrast range, skew residual, speckle density,
scanner make and model from XMP or producer string.

\subsection{3.2 Primary output: the grayscale/bilevel
split}\label{primary-output-the-grayscalebilevel-split}

{\def\LTcaptype{none} % do not increment counter
\begin{longtable}[]{@{}
  >{\raggedright\arraybackslash}p{(\linewidth - 4\tabcolsep) * \real{0.3333}}
  >{\raggedright\arraybackslash}p{(\linewidth - 4\tabcolsep) * \real{0.3333}}
  >{\raggedright\arraybackslash}p{(\linewidth - 4\tabcolsep) * \real{0.3333}}@{}}
\toprule\noalign{}
\begin{minipage}[b]{\linewidth}\raggedright
Bucket
\end{minipage} & \begin{minipage}[b]{\linewidth}\raggedright
Definition
\end{minipage} & \begin{minipage}[b]{\linewidth}\raggedright
Handling
\end{minipage} \\
\midrule\noalign{}
\endhead
\bottomrule\noalign{}
\endlastfoot
\textbf{G} & 8-bit grayscale & Full stroke-morphology detection
available; multi-binarization ensemble applies \\
\textbf{B1} & 1-bit bilevel, lossless (CCITT G4) & Stroke intensity
discarded; morphology degraded but workable \\
\textbf{B2} & 1-bit bilevel, lossy (JBIG2) & As B1, plus
digit-substitution hazard (§3.3) \\
\end{longtable}
}

Bucket G is materially better input than B1 or B2 and should be
preserved as grayscale through the pipeline rather than binarized early.
Binarizing a grayscale page and discarding the original throws away the
information that stroke detection needs, and it is a one-way operation.

\subsection{3.3 JBIG2 hazard}\label{jbig2-hazard}

JBIG2 lossy compression substitutes visually similar glyph patches
across a document and can alter digits with no visible artefact. This
risk is elevated in an engagement when the corpus includes
bilevel-eligible pages.

All numerics from JBIG2 pages are marked suspect regardless of OCR
confidence. Those strata are escalated in QA sampling. If a material
share of the corpus is JBIG2 and any better-quality source exists,
re-deriving those pages is the single highest-return remediation
available.

\subsection{3.4 Image quality scoring and planned remediation
routing}\label{image-quality-scoring-and-planned-remediation-routing}

The implemented profiler records the measurable PDF/image
characteristics and scan branch. The following calibrated quality score
and remediation router is a target integration that requires the client
golden set:

{\def\LTcaptype{none} % do not increment counter
\begin{longtable}[]{@{}
  >{\raggedright\arraybackslash}p{(\linewidth - 2\tabcolsep) * \real{0.5000}}
  >{\raggedright\arraybackslash}p{(\linewidth - 2\tabcolsep) * \real{0.5000}}@{}}
\toprule\noalign{}
\begin{minipage}[b]{\linewidth}\raggedright
Score
\end{minipage} & \begin{minipage}[b]{\linewidth}\raggedright
Path
\end{minipage} \\
\midrule\noalign{}
\endhead
\bottomrule\noalign{}
\endlastfoot
Good & Standard pipeline \\
Marginal & Enhancement pass --- deskew, despeckle, contrast
normalization, and for grayscale, multi-binarization ensemble \\
Poor & Enhancement plus mandatory QA sampling inclusion regardless of
stratum quotas \\
\end{longtable}
}

The implemented \textbf{multi-binarization ensemble} (grayscale pages
only) produces three binarized variants using global thresholding,
adaptive local thresholding, and a stroke-preserving method. OCR each
variant independently. Each variant becomes an additional consensus
voter. On degraded scans this recovers characters that any single
threshold loses, and it costs only compute.

\textbf{Upscaling} for pages below 300 DPI remains planned before any
calibrated handwriting-region processing. Sub-300 DPI is where
handwritten digit recognition fails hardest.

\subsection{3.5 Stamp recovery}\label{stamp-recovery}

Colour separation would have found stamps trivially. The planned
compensating control is template matching against a client-approved
library of recurring ``PAID'', received-date, and approval stamps. It is
not implemented in release 1.0 and must be calibrated before it can
affect payment status.

This matters specifically for payment matching. A ``PAID'' stamp with a
handwritten date is frequently the only evidence that an invoice was
settled, and losing it inflates the apparent open-receivables balance in
§9.

\subsection{3.6 Decision gate}\label{decision-gate}

Outputs: bucket counts, quality distribution, JBIG2 page count, DPI
distribution, and a costed remediation recommendation if better-quality
masters exist.

\begin{center}\rule{0.5\linewidth}{0.5pt}\end{center}

\section{4. Phase 1 --- Ingestion}\label{phase-1-ingestion}

\subsection{4.1 Retain, burst, and prepare
variants}\label{retain-burst-and-prepare-variants}

The implemented intake retains a byte-verified source PDF and creates
one immutable PDF master per page. The implemented preparation ensemble
then creates a grayscale master, contrast-enhanced sibling, and three
complementary binarized siblings without changing any PDF master.
Scan-specific upscaling, deskew, despeckle, and auto-rotation remain
calibrated integration work and must create retained variants rather
than replace evidence.

\subsection{4.2 Text-layer detection}\label{text-layer-detection}

\texttt{pdftotext} is retained per page. In the implemented optional
OpenAI lane, rule-classified pages may use native text and all other
pages use their one-page PDF. Routing for any other OCR provider
requires its own approved adapter and golden-set test.

\subsection{4.3 Document reassembly}\label{document-reassembly}

The implemented control proposes a group only from a complete
consecutive \texttt{PAGE\ X\ OF\ Y} sequence with an agreeing type and
identifier. Its explicit broad mode proposes, but does not order,
same-type/unique-ID candidates. The target signal hierarchy is:

\begin{enumerate}
\def\labelenumi{\arabic{enumi}.}
\tightlist
\item
  OCR-recovered ``Page X of Y'' markers
\item
  Invoice, BOL, AWB, \textbf{ACK, or job number} continuity across
  consecutive pages
\item
  Header-page classifier
\item
  Blank-page or separator-sheet detection
\end{enumerate}

Signal 2 is stronger when attribution keys are extracted as first-class
fields, ACK and job numbers become reassembly evidence as well as
attribution evidence.

Pages that cannot be confidently assigned go to an unassigned pool and
are worked manually. Never force-fit.

\subsection{4.4 OCR engine integration
target}\label{ocr-engine-integration-target}

Release 1.0 includes the optional OpenAI proposal voter, vendor-neutral
OCR/HTR contracts, and a separately enabled Google Document AI
OCR/table-evidence adapter. Document AI retains raw page response and
table/cell evidence but makes no canonical decision; it is a conditional
buddy-audit/reconciliation input. The target production configuration
uses independently tested readers after client selection and golden-set
benchmarking. Bilevel input raises single-engine error rates enough that
a third independent reading should be routine, not a last-resort
tiebreaker.

When provider adapters and the retained variant ensemble are
independently calibrated, one page may contribute multiple readings.
Variants from one engine are correlated evidence and must not be
misrepresented as independent vendors.

\subsection{4.5 Deduplication}\label{deduplication}

Release 1.0 records exact retained-artifact duplicates. Perceptual-image
and normalized-text near-duplicate matching remain a calibrated
proposal-only extension; candidates may never be silently deleted.

\subsection{4.6 Provenance}\label{provenance}

Every record carries source file, page range, engine, engine version,
binarization variant, run timestamp, per-field confidence, and quality
score.

\begin{center}\rule{0.5\linewidth}{0.5pt}\end{center}

\section{5. Phase 2 --- Classification}\label{phase-2-classification}

Two levels: document type, then party role.

Deterministic rules on form numbers, letterhead tokens, and header
strings first; measure against the calibration sample; train a
classifier only if rules plateau below target.

\textbf{Acknowledgements are classified first and processed first} where
they exist as a document class, because they populate the attribution
reference table that Phase 3A depends on.

Below-threshold classifications route to human review. A misclassified
document produces a structurally wrong extraction, not a merely
inaccurate one.

\begin{center}\rule{0.5\linewidth}{0.5pt}\end{center}

\section{6. Phase 3 --- Extraction}\label{phase-3-extraction}

\subsection{6.1 Field schema}\label{field-schema}

Per document type. Header fields are cheap; line items are 3--5× the
effort and are the primary cost driver.

\textbf{Two new mandatory header fields on every financial document
type:} \texttt{ack\_\allowbreak{}number} and
\texttt{job\_\allowbreak{}number}. These are extracted with the same
rigour as monetary fields, because §8 makes them load-bearing.

\subsection{6.2 Independent extraction and
consensus}\label{independent-extraction-and-consensus}

{\def\LTcaptype{none} % do not increment counter
\begin{longtable}[]{@{}
  >{\raggedright\arraybackslash}p{(\linewidth - 4\tabcolsep) * \real{0.3333}}
  >{\raggedright\arraybackslash}p{(\linewidth - 4\tabcolsep) * \real{0.3333}}
  >{\raggedright\arraybackslash}p{(\linewidth - 4\tabcolsep) * \real{0.3333}}@{}}
\toprule\noalign{}
\begin{minipage}[b]{\linewidth}\raggedright
Condition
\end{minipage} & \begin{minipage}[b]{\linewidth}\raggedright
Disposition
\end{minipage} & \begin{minipage}[b]{\linewidth}\raggedright
Flag
\end{minipage} \\
\midrule\noalign{}
\endhead
\bottomrule\noalign{}
\endlastfoot
Independent primary and secondary readings agree & Consensus proposal;
applicable deterministic controls still run &
\texttt{consensus\_\allowbreak{}confirmed} \\
Independent readings differ or a provider fails & Explicit exception;
never silently accepted & \texttt{no\_\allowbreak{}consensus} \\
A later independent corroboration is available & Retain it as additional
evidence; it does not erase prior disagreement &
\texttt{consensus\_\allowbreak{}corroborated} \\
\end{longtable}
}

Where the multi-binarization ensemble applies, variant readings are
additional evidence within one provider's result rather than independent
proof. A field where binarization variants disagree is a signal that the
page is degraded at that location, and it is flagged even if the
providers agree overall.

Provider budgets are set per engagement. Independent evidence is
preserved because a wrong figure can propagate through a long-lived
financial history.

\subsection{6.3 Arithmetic self-proof}\label{arithmetic-self-proof}

Line extensions reconcile to line totals; lines sum to subtotal;
subtotal + tax + freight + accessorials = stated total, to the cent.

A document failing its own arithmetic is an exception regardless of
confidence. It can also be the primary means of discovering handwritten
overrides that detection missed (§1.2).

\subsection{6.4 Additional validation}\label{additional-validation}

Dates within the operating window; currency and UOM against controlled
vocabularies; invoice numbers against detected per-vendor format
patterns; \textbf{ACK and job numbers against the format specification
from §2.3}; tax rates within plausible ranges.

Format validation on attribution keys is the cheapest error detection in
the project. A five-character ACK number in a corpus where every other
ACK is six characters is an OCR error, and the check costs nothing.

Raw addresses are retained as evidence and locally split into CRM
components: address line 1, address line 2, city, state or region,
postal code, and ISO country code. An approved run may then explicitly
call Google Address Validation with a non-secret API key and a maximum
of 5,000 provider requests. It records a separate standardized-address,
deliverability, and geocode proposal only when the verdict is complete
and premise-level without component uncertainty. Provider failures,
including unsupported country coverage, caps, incomplete verdicts, and
conflicts stay in the final-review queue; no source value is
overwritten. Address Validation itself supplies the geocode evidence, so
this path does not call the separate Geocoding API.

\subsection{6.5 Exception policy}\label{exception-policy}

Sampling governs the clean population only. Every record failing
consensus, arithmetic, validation, or reconciliation is worked by a
human, 100\% of the time. Expect this to run at the upper end of the
typical 5--15\% band given bilevel input.

\subsection{6.6 Entity resolution}\label{entity-resolution}

Party names and addresses normalized and fuzzy-matched into a
deduplicated master with defined surviving-record rules. ``ACME Corp'',
``Acme Corporation'', and ``ACME CORP.'' must collapse to one account or
the CRM is worthless on delivery.

Every merge decision is logged and reversible. Ambiguous clusters are
presented for adjudication, not guessed at.

\begin{center}\rule{0.5\linewidth}{0.5pt}\end{center}

\section{7. Phase 3H --- Handwriting Capture
Track}\label{phase-3h-handwriting-capture-track}

\subsection{7.1 Scope decision and its
boundary}\label{scope-decision-and-its-boundary}

The client has directed that handwriting may be ignored. That direction
is accepted for \textbf{recognition} --- full handwritten text
recognition with 3-of-3 consensus is descoped, along with its manual
review burden.

It cannot be accepted for \textbf{capture}, for two specific reasons:

\textbf{First, the attribution requirement.} §2.3 asks where ACK and job
numbers physically appear. On shipping paperwork these are very often
handwritten --- written on at the point of order entry or at receiving.
If handwriting is discarded entirely and ACK numbers are among the
discarded content, the attribution requirement in §8 cannot be met at
all. The two instructions are in direct tension, and attribution is the
harder requirement.

\textbf{Second, arithmetic integrity.} Handwritten quantity and price
corrections change what a document actually says. Ignoring them does not
make the printed figure correct; it makes it wrong and unflagged.
Arithmetic proof will surface these regardless (§1.2), so the work of
dealing with them is not avoided by descoping --- only the ability to
resolve them cheaply is lost.

\subsubsection{The resulting scope}\label{the-resulting-scope}

{\def\LTcaptype{none} % do not increment counter
\begin{longtable}[]{@{}
  >{\raggedright\arraybackslash}p{(\linewidth - 2\tabcolsep) * \real{0.5000}}
  >{\raggedright\arraybackslash}p{(\linewidth - 2\tabcolsep) * \real{0.5000}}@{}}
\toprule\noalign{}
\begin{minipage}[b]{\linewidth}\raggedright
Activity
\end{minipage} & \begin{minipage}[b]{\linewidth}\raggedright
Status
\end{minipage} \\
\midrule\noalign{}
\endhead
\bottomrule\noalign{}
\endlastfoot
Detect handwritten regions & \textbf{Retained} --- cheap, and required
for the flag \\
Crop and store region images & \textbf{Retained} --- the artefact the
CRM links to \\
Read ACK / job numbers from handwriting & \textbf{Retained} --- mandated
by §8 \\
Read quantities/amounts where arithmetic fails & \textbf{Retained} ---
triggered by failure, not by default \\
Read cheque numbers and payment marks & \textbf{Retained} --- needed for
§9 aging closure \\
Read all other handwritten content & \textbf{Descoped} --- captured as
an image, flagged, not transcribed \\
3-of-3 consensus on all handwritten numerics & \textbf{Narrowed} ---
applies only to the retained categories \\
Manual review of every handwritten amendment & \textbf{Narrowed} ---
applies only where a dollar figure or attribution key is affected \\
\end{longtable}
}

This preserves the client's cost saving --- the bulk of handwriting is
delivery notes, initials, and marginalia that genuinely does not need
transcription --- while protecting the two things that would otherwise
break.

\subsection{7.2 Detection without colour: target
architecture}\label{detection-without-colour-target-architecture}

Three signals in combination, since ink separation is unavailable.

\textbf{Stroke morphology} --- planned local stroke-width variance,
baseline instability, slant irregularity, curvature entropy. Primary
detector. Degraded on bilevel relative to grayscale, which is why the
§3.2 bucket split matters.

\textbf{Layout delta} --- planned registration against a learned,
client-approved form template; ink outside the printed field grid or
overlapping printed glyphs is annotation. Strong on standardized forms,
weak on free-format documents.

\textbf{Full-page vision-language model proposal} --- the implemented
optional PDF-mode OpenAI lane may enumerate likely handwritten regions
with bounding boxes. It is proposal-only and cannot substitute for a
calibrated detector or independent HTR.

Detection recall will be materially below what colour scans would give.
Detector output sets processing priority, never eligibility.

\subsection{7.3 Capture and CRM linkage
contract}\label{capture-and-crm-linkage-contract}

The canonical contract supports a record for every detected region
regardless of whether it is transcribed:

\begin{verbatim}
handwriting_region:
  document_id, page, bounding_box, region_image_path,
  detector, detector_confidence,
  annotation_type_guess, transcribed (bool),
  transcribed_text (nullable), recognition_agreement (nullable)
\end{verbatim}

The target CRM adapter stores and links the cropped region image. The
practical effect: a user viewing an invoice sees a flag that the
document carries handwriting, and can click through to the actual
cropped image of it without anyone having transcribed it. This satisfies
the client's instruction to ``capture it somehow and log it'' at
near-zero marginal cost, and it means the decision to transcribe can be
revisited later without reprocessing the corpus.

\subsection{7.4 Targeted recognition
triggers}\label{targeted-recognition-triggers}

Transcription runs only when triggered:

\begin{enumerate}
\def\labelenumi{\arabic{enumi}.}
\tightlist
\item
  Region classified as likely containing an ACK or job number, by
  position and character pattern
\item
  Document failed arithmetic self-proof and a handwritten region
  overlaps a quantity, price, or total field
\item
  Region overlaps a payment block, or a ``PAID'' stamp was matched
  nearby
\item
  Region falls within a QA sample selection
\end{enumerate}

The target integration uses dedicated independent HTR engines in
parallel, full-page context rather than isolated crops, 3-of-3 agreement
required for numerics. The strictness is retained where it is applied;
only the scope of application is narrowed.

\subsection{7.5 Precedence rule}\label{precedence-rule}

Unchanged and non-negotiable:

\begin{quote}
\textbf{Printed values are the baseline. Handwritten values are
amendments carrying a later effective time. Both are stored.}
\end{quote}

The canonical layer exposes the amended value as current and the printed
value as \texttt{original\_\allowbreak{}value}, with
\texttt{amendment\_\allowbreak{}source = 'handwritten'}. Auditable and
reversible.

\subsection{7.6 Signatures}\label{signatures}

Never transcribed as data. Recorded as \texttt{signed} (boolean), region
bounding box, and approver identity only where separately legible from a
printed name line. Reading cursive signatures as names produces
confident nonsense that contaminates entity resolution.

\begin{center}\rule{0.5\linewidth}{0.5pt}\end{center}

\section{8. Phase 3A --- Attribution}\label{phase-3a-attribution}

New phase. Every dollar in the dataset ties to an ACK or job/project
number, and every sale carries an originating city. This is a hard gate,
not an enrichment.

\subsection{8.1 Why this is a separate
phase}\label{why-this-is-a-separate-phase}

Attribution is not a field extraction problem. Roughly half the
documents in a corpus of this kind will not state the attribution key
directly --- a carrier freight bill carries a PRO number and a BOL
reference, not a job number. The key has to be \textbf{derived through
the document graph}, which requires extraction and entity resolution to
be complete first.

Treating it as just another extracted field is the most common way this
requirement fails. It has its own phase, its own resolution hierarchy,
its own coverage metric, and its own gate.

\subsection{8.2 Attribution reference
table}\label{attribution-reference-table}

Built first, before any resolution is attempted. Sources in order of
authority:

\begin{enumerate}
\def\labelenumi{\arabic{enumi}.}
\tightlist
\item
  Client's order-entry or quoting system export, if one exists --- a
  digital list of ACK numbers with dates, customers, and amounts. This
  is by far the best input and §2.4 asks for it specifically.
\item
  Acknowledgement documents in the corpus, extracted in Phase 3.
\item
  ACK and job numbers appearing on any other document type.
\end{enumerate}

The reference table gives every known ACK and job number a validated
format, a date range, a customer, and an expected value. Every
subsequent resolution is checked against it, which converts attribution
from guesswork into lookup.

\subsection{8.3 Resolution hierarchy}\label{resolution-hierarchy}

Applied in order. The first method that resolves wins, and \textbf{the
method is recorded on the record}.

{\def\LTcaptype{none} % do not increment counter
\begin{longtable}[]{@{}
  >{\raggedright\arraybackslash}p{(\linewidth - 6\tabcolsep) * \real{0.2500}}
  >{\raggedright\arraybackslash}p{(\linewidth - 6\tabcolsep) * \real{0.2500}}
  >{\raggedright\arraybackslash}p{(\linewidth - 6\tabcolsep) * \real{0.2500}}
  >{\raggedright\arraybackslash}p{(\linewidth - 6\tabcolsep) * \real{0.2500}}@{}}
\toprule\noalign{}
\begin{minipage}[b]{\linewidth}\raggedright
Rank
\end{minipage} & \begin{minipage}[b]{\linewidth}\raggedright
Method
\end{minipage} & \begin{minipage}[b]{\linewidth}\raggedright
Confidence
\end{minipage} & \begin{minipage}[b]{\linewidth}\raggedright
Notes
\end{minipage} \\
\midrule\noalign{}
\endhead
\bottomrule\noalign{}
\endlastfoot
1 & Explicit printed ACK or job number on the document & Highest &
Validated against reference table \\
2 & Explicit handwritten ACK or job number & High & Requires 3-of-3
recognition per §7.4 \\
3 & Purchase order number cross-reference & High & PO → ACK via
reference table \\
4 & Invoice-to-shipment linkage & Medium & Shipment already attributed;
invoice inherits \\
5 & Payment-to-invoice linkage & Medium & Invoice already attributed;
payment inherits \\
6 & Customer + date-window + amount match against reference table &
Medium-low & Requires a unique match; ambiguous matches are not
resolved \\
7 & Email evidence & Low & Phase 7 only. Corroborating, never
overriding \\
8 & Manual assignment & Recorded & Assigner and date logged \\
--- & \textbf{Unattributable} & --- & Explicit disposition with a stated
reason and dollar value \\
\end{longtable}
}

Ranks 4 and 5 are inheritance rules and they are what make the
requirement achievable. A carrier freight bill that states no job number
inherits one from the shipment it bills for, which inherited it from the
invoice, which stated it explicitly. The chain is recorded so the
inference is auditable.

\subsection{8.4 Sales origin city}\label{sales-origin-city}

The client requires knowing what city a sale came from. \textbf{This
phrase has at least four defensible meanings and they produce different
numbers.} Settle it in Phase 0 rather than discovering the disagreement
after the analysis is presented.

{\def\LTcaptype{none} % do not increment counter
\begin{longtable}[]{@{}
  >{\raggedright\arraybackslash}p{(\linewidth - 4\tabcolsep) * \real{0.3333}}
  >{\raggedright\arraybackslash}p{(\linewidth - 4\tabcolsep) * \real{0.3333}}
  >{\raggedright\arraybackslash}p{(\linewidth - 4\tabcolsep) * \real{0.3333}}@{}}
\toprule\noalign{}
\begin{minipage}[b]{\linewidth}\raggedright
Interpretation
\end{minipage} & \begin{minipage}[b]{\linewidth}\raggedright
Meaning
\end{minipage} & \begin{minipage}[b]{\linewidth}\raggedright
Typical source
\end{minipage} \\
\midrule\noalign{}
\endhead
\bottomrule\noalign{}
\endlastfoot
Selling location & Which of the client's offices or branches booked the
sale & Letterhead, issuing office, ACK prefix, salesperson \\
Ship-from origin & Where the goods physically departed & BOL origin,
warehouse address \\
Customer location & Where the customer is based & Bill-to address \\
Delivery destination & Where the goods went & Ship-to address, POD \\
\end{longtable}
}

In most businesses asking this question, the intended meaning is
\textbf{selling location} --- a territory or branch performance
question. It is also the hardest of the four to source, because it is
frequently not stated on the document at all and has to be inferred from
letterhead, ACK number prefix, or salesperson assignment.

\subsubsection{Resolution hierarchy for selling
location}\label{resolution-hierarchy-for-selling-location}

\begin{enumerate}
\def\labelenumi{\arabic{enumi}.}
\tightlist
\item
  Explicit branch or office field on the document
\item
  Letterhead or footer address on the issuing document --- recoverable
  by template matching against a small library of the client's
  letterhead variants
\item
  ACK number prefix, where the numbering scheme encodes location
  (common; check in §2.3)
\item
  Salesperson or rep name, mapped to a territory table
\item
  Customer's assigned territory from the reference table
\item
  Email evidence (Phase 7)
\item
  Unresolved --- explicit, with dollar value
\end{enumerate}

All four interpretations are captured where the data supports it, stored
as separate fields. This costs little and means the client can change
their mind about which one they meant without reprocessing.

\subsection{8.5 Attribution coverage
gate}\label{attribution-coverage-gate}

The exit condition for this phase.

{\def\LTcaptype{none} % do not increment counter
\begin{longtable}[]{@{}
  >{\raggedright\arraybackslash}p{(\linewidth - 2\tabcolsep) * \real{0.5000}}
  >{\raggedright\arraybackslash}p{(\linewidth - 2\tabcolsep) * \real{0.5000}}@{}}
\toprule\noalign{}
\begin{minipage}[b]{\linewidth}\raggedright
Metric
\end{minipage} & \begin{minipage}[b]{\linewidth}\raggedright
Requirement
\end{minipage} \\
\midrule\noalign{}
\endhead
\bottomrule\noalign{}
\endlastfoot
Dollars attributed to ACK or job number & 100\% attributed \textbf{or}
explicitly registered unattributable \\
Attribution by method & Reported --- the mix of ranks 1--8 is a quality
indicator in itself \\
Dollars with resolved selling location & Reported, with unresolved
quantified \\
Unattributable register & Every entry has a reason code and a dollar
value \\
\end{longtable}
}

An attribution rate of ``97\%'' is not a passing result unless the
remaining 3\% is enumerated with reasons and a dollar figure. The gate
is not high attribution; it is that \textbf{no dollar is unaccounted for
silently.}

Reason codes for legitimate unattributables --- internal transfer,
rebate, rebilled freight, pre-system-implementation, credit against a
closed job --- are agreed in Phase 0 so they are dispositions rather
than failures.

\begin{center}\rule{0.5\linewidth}{0.5pt}\end{center}

\section{9. Phase 4 --- Completeness and Data Quality
Gate}\label{phase-4-completeness-and-data-quality-gate}

Completeness is a harder requirement than accuracy. A missing document
is invisible; a wrong one is at least detectable.

\subsection{9.1 Financial event stream}\label{financial-event-stream}

All documents normalize to a ledger-shaped stream:

\begin{verbatim}
event_id, event_type, party_id, amount, currency, effective_date,
source_document_id, ack_number, job_number, selling_location
\end{verbatim}

\texttt{event\_\allowbreak{}type in \{purchase, payment, credit, adjustment, accrual\}}

Attribution keys travel on the event, not just on the source document,
so every downstream aggregation can be cut by job without a join back
through the document graph.

Non-financial documents register with a null amount so the shipment
lifecycle stays complete.

\subsection{9.2 Matching}\label{matching}

\textbf{Three-way match:} PO \textless-\textgreater{} receipt (packing
list or POD) \textless-\textgreater{} invoice.

\textbf{Open-item matching:} invoice \textless-\textgreater{} payment
\textless-\textgreater{} remittance advice, supporting partial payments
and one-payment-to-many-invoices.

\textbf{ACK-to-invoice match:} every ACK in the reference table should
resolve to one or more invoices. ACKs with no invoice are either open
jobs or missing documents --- the distinction matters and is reported.

Every match carries a method and confidence.

\subsection{9.3 Completeness controls}\label{completeness-controls}

\begin{enumerate}
\def\labelenumi{\arabic{enumi}.}
\tightlist
\item
  \textbf{Sequence continuity} --- invoice numbering per vendor; all
  gaps enumerated. When acknowledgements are in scope, apply the control
  to \textbf{ACK number sequences}, which are typically client-issued
  and therefore far more reliably sequential than vendor invoice
  numbers. A gap in the client's own ACK sequence is strong evidence of
  a missing job.
\item
  \textbf{Calendar continuity} --- document counts and amounts by vendor
  by month.
\item
  \textbf{Period reconciliation} --- extracted totals vs GL by period
  and vendor. Variance reported, never plugged.
\item
  \textbf{Bank/cash reconciliation} --- payment events vs bank statement
  totals.
\item
  \textbf{Aging closure} --- every invoice terminates in a payment,
  credit, write-off, or explicit open status.
\item
  \textbf{Attribution closure} --- every dollar attributed or explicitly
  registered (§8.5).
\end{enumerate}

\subsection{9.4 Completeness Report}\label{completeness-report}

A standing deliverable regenerated on every run: total documents and
value; unmatched count and value by cause; sequence gaps including ACK
gaps; calendar gaps; GL variance; \textbf{attribution coverage and the
unattributable register}; exception queue volume and clearance.

Every analytic output carries these figures on its face.

\subsection{9.5 Gate}\label{gate}

Does not proceed to analytics until the one-sided MUS bound is completed
and within tolerance (or the affected population was fully reworked),
the Completeness Report has \texttt{gate\_\allowbreak{}status: clear},
separately reviewed attribute results are accepted, and the §8.5
attribution gate is met.

\begin{center}\rule{0.5\linewidth}{0.5pt}\end{center}

\section{10. Phase 4Q --- Sampling QA}\label{phase-4q-sampling-qa}

Two schemes in parallel, both drawing only from records marked
\texttt{auto\_\allowbreak{}accepted} whose arithmetic status is
explicitly \texttt{proved} or \texttt{not\_\allowbreak{}applicable}.

\subsection{10.1 Monetary unit sampling}\label{monetary-unit-sampling}

Selection probability proportional to document value. Everything above
the materiality threshold, plus any item larger than the provisional
systematic interval, is examined with certainty. Produces a one-sided
projected overstatement and dollar-denominated 95\% upper overstatement
bound. Observed understatement and completeness are reported separately.

\subsection{10.2 Stratified attribute
sampling}\label{stratified-attribute-sampling}

The checked-in script creates a reproducible document-selection plan by
document type × year × processing branch ×
\texttt{has\_\allowbreak{}handwriting}; JBIG2 status is a separate
oversampling trigger. It reports a zero-finding rule-of-three planning
bound per stratum, not observed field accuracy. Vendor rank,
image-quality/scan-bucket, period-edge, and ACK/job field strata require
an engagement-specific extension and validation before they can be
reported.

\textbf{Oversampling:}

{\def\LTcaptype{none} % do not increment counter
\begin{longtable}[]{@{}
  >{\raggedright\arraybackslash}p{(\linewidth - 2\tabcolsep) * \real{0.5000}}
  >{\raggedright\arraybackslash}p{(\linewidth - 2\tabcolsep) * \real{0.5000}}@{}}
\toprule\noalign{}
\begin{minipage}[b]{\linewidth}\raggedright
Stratum
\end{minipage} & \begin{minipage}[b]{\linewidth}\raggedright
Multiplier
\end{minipage} \\
\midrule\noalign{}
\endhead
\bottomrule\noalign{}
\endlastfoot
Handwriting-bearing Branch A & 4× \\
Handwriting-bearing Branch B & 9× \\
Handwriting-bearing B-rescan & 4× \\
JBIG2-suspect stratum & at least 5× \\
\end{longtable}
}

Attribution keys are a recommended engagement-specific oversample
because a wrong ACK number does not produce a wrong total. The current
generic planner does not derive a field-level ACK/job stratum, so an
operator must not claim it did.

\subsection{10.3 Sizing}\label{sizing}

Rule of three: zero errors in \emph{n} gives an upper 95\% bound of
approximately 3/\emph{n}.

{\def\LTcaptype{none} % do not increment counter
\begin{longtable}[]{@{}
  >{\raggedright\arraybackslash}p{(\linewidth - 2\tabcolsep) * \real{0.5000}}
  >{\raggedright\arraybackslash}p{(\linewidth - 2\tabcolsep) * \real{0.5000}}@{}}
\toprule\noalign{}
\begin{minipage}[b]{\linewidth}\raggedright
n
\end{minipage} & \begin{minipage}[b]{\linewidth}\raggedright
Upper bound
\end{minipage} \\
\midrule\noalign{}
\endhead
\bottomrule\noalign{}
\endlastfoot
60 & \textasciitilde5\% \\
300 & \textasciitilde1\% \\
600 & \textasciitilde0.5\% \\
\end{longtable}
}

Sizing derives from tolerable error rate, not corpus size.

\subsection{10.4 Escalation}\label{escalation}

If the completed MUS upper overstatement bound exceeds tolerance, the
affected population goes to 100\% review and is re-sampled after
remediation. Invalid, duplicate, off-sample, or
confidence-factor-table-exhausting findings block the calculation
instead of producing a bound.

\subsection{10.5 Golden set}\label{golden-set}

500--1,000 hand-labelled pages spanning every document type, vendor
format, quality band, both bilevel and grayscale buckets, and heavy
handwriting representation. Permanent regression infrastructure; every
pipeline change is benchmarked against it before deployment.

\begin{center}\rule{0.5\linewidth}{0.5pt}\end{center}

\section{11. Phase 5 --- Analytics}\label{phase-5-analytics}

Scoped against the Phase 0 decision list.

\subsection{11.1 Job and ACK analytics}\label{job-and-ack-analytics}

When configured under §8, this is a primary analytic frame.

\begin{itemize}
\tightlist
\item
  Revenue, cost, and margin by job or ACK
\item
  Job profitability distribution; loss-making job identification
\item
  Cost composition by job --- materials, freight, accessorials, duty
\item
  Job cycle time: ACK date → ship date → invoice date → payment date
\item
  Open ACKs with no invoice; invoiced jobs with no payment
\item
  Revenue by job over time; job size distribution
\end{itemize}

\subsection{11.2 Geographic analytics}\label{geographic-analytics}

\begin{itemize}
\tightlist
\item
  Revenue by selling location, by period
\item
  Growth and decline by location
\item
  Customer concentration within each location
\item
  Lane and destination mix by originating location
\item
  Margin by location
\end{itemize}

Every geographic figure carries its resolution-method mix. A city total
that is 60\% inferred from ACK prefix and 40\% explicitly stated is a
different quality of number than one that is fully explicit, and the
reader needs to know which they are looking at.

\subsection{11.3 Customer analytics}\label{customer-analytics}

Revenue by account and period; concentration; first-order date; RFM;
churn signal; cohort retention; product mix.

Churn findings are cross-checked against the calendar-gap report before
presentation. A customer who appears to have stopped ordering may simply
have documents in an unscanned batch --- the most common false finding
in backfile analytics.

\subsection{11.4 Vendor and carrier
analytics}\label{vendor-and-carrier-analytics}

Spend by carrier; rate variance by lane; on-time delivery from POD
dates; claim and damage rate; vendor concentration; fuel surcharge
ratio.

\subsection{11.5 Cost analytics}\label{cost-analytics}

Freight cost per shipment, per unit weight, per lane;
\textbf{accessorial leakage} --- detention, demurrage, fuel surcharge,
redelivery, storage; landed cost; cost trend normalized for volume.

Accessorial leakage is usually where recoverable money is found. It
requires accessorials itemized by code at extraction; collapsing them
into a single freight figure makes the analysis impossible without
re-extraction.

\subsection{11.6 Financial control
analytics}\label{financial-control-analytics}

Duplicate invoice detection across the full corpus; price-vs-contract
variance; payment terms adherence; DSO/DPO; early-payment discount
capture; unapplied cash; overbilling via three-way match variance.

\subsection{11.7 Data quality analytics}\label{data-quality-analytics}

A permanent operational dashboard: extraction accuracy by field, type,
and quality band; handwriting incidence and capture rate;
\textbf{attribution coverage by method}; exception volume and clearance;
completeness metrics; entity resolution status.

\begin{center}\rule{0.5\linewidth}{0.5pt}\end{center}

\section{12. Phase 6 --- Data Model and CRM
Handoff}\label{phase-6-data-model-and-crm-handoff}

\subsection{12.1 Layers}\label{layers}

\textbf{Raw} (immutable engine output, retained permanently) →
\textbf{staged} (typed, validated) → \textbf{canonical} (deduplicated,
entity-resolved, amendments applied) → \textbf{CRM load}
(target-shaped).

Any downstream layer rebuilds from raw without re-OCR.

\subsection{12.2 Core entities}\label{core-entities}

\texttt{party}, \texttt{party\_\allowbreak{}role}, \texttt{address},
\texttt{contact}, \texttt{item}, \texttt{carrier}, \texttt{lane},
\texttt{shipment}, \texttt{document},
\texttt{invoice\_\allowbreak{}header},
\texttt{invoice\_\allowbreak{}line}, \texttt{payment},
\texttt{payment\_\allowbreak{}application}, \texttt{amendment},
\texttt{exception\_\allowbreak{}event}

\textbf{Implemented control:}

{\def\LTcaptype{none} % do not increment counter
\begin{longtable}[]{@{}
  >{\raggedright\arraybackslash}p{(\linewidth - 2\tabcolsep) * \real{0.5000}}
  >{\raggedright\arraybackslash}p{(\linewidth - 2\tabcolsep) * \real{0.5000}}@{}}
\toprule\noalign{}
\begin{minipage}[b]{\linewidth}\raggedright
Entity
\end{minipage} & \begin{minipage}[b]{\linewidth}\raggedright
Purpose
\end{minipage} \\
\midrule\noalign{}
\endhead
\bottomrule\noalign{}
\endlastfoot
\texttt{acknowledgement} & The ACK master. Number, date, customer,
selling location, value, status \\
\texttt{job} & Job or project master where distinct from ACK \\
\texttt{attribution} & Links any financial row to an ACK or job, with
method, confidence, and evidence chain \\
\texttt{selling\_\allowbreak{}location} & Branch/office master with
territory mapping \\
\texttt{handwriting\_\allowbreak{}region} & Detected annotation regions,
cropped images, transcription status \\
\texttt{email\_\allowbreak{}message} & Phase 7 only \\
\texttt{email\_\allowbreak{}evidence} & Phase 7 only --- links an email
to a document, attribution, or party \\
\end{longtable}
}

\subsection{12.3 Common control columns}\label{common-control-columns}

Every batch-managed business row carries its idempotency key,
\texttt{batch\_\allowbreak{}id}, and
\texttt{review\_\allowbreak{}status}. Evidence-bearing rows also carry a
direct source-document key or explicit decision-evidence link. Static
vocabulary and pure association rows retain their own stable
natural/composite keys without inventing page provenance.

Applicable transaction fields include:

\begin{itemize}
\tightlist
\item
  Identity: stable primary/idempotency key and
  \texttt{natural\_\allowbreak{}key} where applicable.
\item
  Source proof: \texttt{source\_\allowbreak{}document\_\allowbreak{}id}
  and \texttt{source\_\allowbreak{}page\_\allowbreak{}range}.
\item
  Control: \texttt{source\_\allowbreak{}confidence},
  \texttt{review\_\allowbreak{}status}, and
  \texttt{has\_\allowbreak{}handwriting}.
\item
  Value lineage: \texttt{original\_\allowbreak{}value},
  \texttt{amended\_\allowbreak{}value}, and
  \texttt{amendment\_\allowbreak{}source}.
\item
  Quality band: \texttt{image\_\allowbreak{}quality\_\allowbreak{}band}.
\item
  Scan bucket: \texttt{scan\_\allowbreak{}bucket}.
\item
  Processing audit: \texttt{created\_\allowbreak{}at},
  \texttt{updated\_\allowbreak{}at}, and
  \texttt{batch\_\allowbreak{}id}.
\end{itemize}

Financial tables carry \texttt{ack\_\allowbreak{}number} and
\texttt{job\_\allowbreak{}number}. They also carry
\texttt{attribution\_\allowbreak{}method} and
\texttt{selling\_\allowbreak{}location\_\allowbreak{}key}.

\subsection{12.4 Load order}\label{load-order}

\begin{enumerate}
\def\labelenumi{\arabic{enumi}.}
\tightlist
\item
  Reference data --- currency, UOM, charge codes, country, document
  type, \textbf{selling\_location}
\item
  \texttt{document} registry --- before every source-provenance foreign
  key
\item
  \texttt{party} → \texttt{address} → \texttt{contact}
\item
  \texttt{carrier}, \texttt{item}, \texttt{lane}
\item
  \textbf{\texttt{acknowledgement}, \texttt{job}}
\item
  \texttt{shipment}
\item
  \texttt{invoice\_\allowbreak{}header}
\item
  \texttt{invoice\_\allowbreak{}line}
\item
  \texttt{payment}, then \texttt{payment\_\allowbreak{}application}
\item
  \textbf{\texttt{attribution}}
\item
  \textbf{\texttt{handwriting\_\allowbreak{}region}},
  \texttt{amendment}, \texttt{exception\_\allowbreak{}event}
\item
  Phase 7 only: \texttt{email\_\allowbreak{}message},
  \texttt{email\_\allowbreak{}evidence}
\end{enumerate}

ACK and job masters load before any transactional table that references
them. Attribution loads after both sides exist.

\subsection{12.5 Load mechanics}\label{load-mechanics}

Idempotent upsert by natural key. Referential integrity enforced at
load, not assumed --- a failed foreign key halts the load rather than
silently dropping the row. Load manifest per run. Rollback by
\texttt{batch\_\allowbreak{}id}.

\subsection{12.6 Handoff deliverables}\label{handoff-deliverables}

Canonical dataset; ERD and data dictionary; load order specification;
integrity rules; Completeness Report; one-sided MUS overstatement bound
and separately reviewed attribute-sample results; \textbf{attribution
coverage report and unattributable register}; open exception register.

\begin{center}\rule{0.5\linewidth}{0.5pt}\end{center}

\section{13. Phase 7 --- Email Corpus Integration (Optional,
Deferred)}\label{phase-7-email-corpus-integration-optional-deferred}

An email corpus may be incorporated as an ancillary step when separately
authorized. It is designed so that deferring it costs nothing and
executing it requires no rework of Phases 0--6.

\subsection{13.1 Scope mismatch, stated up
front}\label{scope-mismatch-stated-up-front}

An email corpus may cover a shorter or different period than the
document corpus. \textbf{Email can enrich only its measured overlap
window.} Any analytic conclusion drawn from email-derived data applies
to that window and must state the coverage limit. Presenting an
email-enriched metric alongside a full-period trend without noting the
coverage break produces a false discontinuity that looks like a business
change.

\subsection{13.2 What email is good for}\label{what-email-is-good-for}

In descending order of likely value:

\textbf{1. Attachments.} Emails frequently contain a digitally generated
PDF of an invoice that exists in the scan corpus only as a degraded
black-and-white image. Where an attachment matches a scan, retain it as
a separate source with its own hash and provenance. Compare its derived
values to the scan/OCR proposals and record agreement or a cited
amendment/discrepancy; never replace either source. Given that the
historic corpus is bilevel, this is the single most valuable thing in
the mailbox and on its own may justify the phase.

\textbf{2. Attribution gap-filling.} Emails carry ACK numbers, job
numbers, and PO references in subject lines and bodies. For dollars
stranded at rank 8 or unattributable in §8.3, email is often the only
remaining evidence.

\textbf{3. Missing payment evidence.} Remittance advices, payment
confirmations, and ``cheque is in the mail'' threads close aging items
that the document corpus leaves open.

\textbf{4. Selling location evidence.} Signature blocks carry office
addresses. Sender identity maps to territory. This resolves §8.4 rank 6.

\textbf{5. Contact enrichment for the CRM.} Names, titles, phone
numbers, and current email addresses per account --- data the document
corpus barely contains and the CRM genuinely needs.

\textbf{6. Order-to-cash narrative.} Quote → ACK → change orders →
dispute → settlement threads explain anomalies that the documents only
record the outcome of.

\subsection{13.3 What email must not be used
for}\label{what-email-must-not-be-used-for}

\textbf{Email is corroborating evidence. It never overrides a source
document.}

An email saying ``invoice 2231 is for \$5,500'' does not correct an
invoice that says \$5,507.50. It raises a discrepancy for human
resolution. The exception is §13.2 item 1: a PDF \emph{attachment} is a
document, not an email, and is treated as such.

All email-derived values carry
\texttt{source\ =\ \textquotesingle{}email\textquotesingle{}} and are
excluded from the document-based accuracy statement, which is computed
on documents only.

\subsection{13.4 Process}\label{process}

\textbf{Screening first.} Before any content processing: define mailbox
scope, date range, and an exclusion filter for personal correspondence,
HR matters, legal advice, and anything privileged. Screening precedes
ingestion, not follows it. The output of this step is a written scope
agreement.

\textbf{Ingest.} PST, MBOX, EML, or Graph API export. Normalize to a
common structure: message id, thread id, sender, recipients, date,
subject, body, attachments.

\textbf{Thread reconstruction and deduplication.} Reply chains duplicate
body text extensively; without deduplication the corpus inflates
several-fold and reference extraction returns the same fact dozens of
times.

\textbf{Attachment extraction and matching.} Every attachment extracted,
hashed, and matched against the document corpus by content and by
reference number. Matched attachments are promoted into the document
pipeline as higher-fidelity replacements, processed through Phases 2--3
normally, and their extracted values compared against the scan-derived
values. \textbf{Disagreements are findings}, not silent overwrites ---
they measure how much the black-and-white scanning actually cost.

\textbf{Reference extraction.} ACK numbers, job numbers, invoice
numbers, PO numbers, cheque numbers, and amounts, by pattern match
against the §8.2 reference table formats, plus a language-model pass for
references stated in prose.

\textbf{Entity linkage.} Sender and recipient domains and addresses
matched to the resolved party master from §6.6. Contacts extracted for
CRM load.

\textbf{Evidence linkage.} Each extracted reference becomes an
\texttt{email\_\allowbreak{}evidence} row linking a message to a
document, attribution, or party, with the quoted text and a confidence.

\subsection{13.5 Deliverables}\label{deliverables-1}

Email-derived contact list for CRM load; attribution gap-fill report
with before/after coverage; recovered payment evidence with its effect
on aging closure; \textbf{attachment-versus-scan discrepancy report};
revised Completeness Report for the email-covered window.

\subsection{13.6 Why this is deferred rather than
dropped}\label{why-this-is-deferred-rather-than-dropped}

Executing Phase 7 later costs nothing extra provided Phases 0--6 leave
room for it. The room required is: \texttt{source} fields on every
value, an \texttt{attribution} table that accepts new evidence rows
without schema change, and an unattributable register that persists
rather than being cleared at handoff. All three are in the The current
design.

\begin{center}\rule{0.5\linewidth}{0.5pt}\end{center}

\section{14. Build Sequence and
Parallelization}\label{build-sequence-and-parallelization}

{\def\LTcaptype{none} % do not increment counter
\begin{longtable}[]{@{}
  >{\raggedright\arraybackslash}p{(\linewidth - 4\tabcolsep) * \real{0.3333}}
  >{\raggedright\arraybackslash}p{(\linewidth - 4\tabcolsep) * \real{0.3333}}
  >{\raggedright\arraybackslash}p{(\linewidth - 4\tabcolsep) * \real{0.3333}}@{}}
\toprule\noalign{}
\begin{minipage}[b]{\linewidth}\raggedright
Track
\end{minipage} & \begin{minipage}[b]{\linewidth}\raggedright
Dependency
\end{minipage} & \begin{minipage}[b]{\linewidth}\raggedright
May start
\end{minipage} \\
\midrule\noalign{}
\endhead
\bottomrule\noalign{}
\endlastfoot
Scan profiling (0.5) & Corpus access & Immediately \\
Ingestion (1) & Phase 0 taxonomy & After Phase 0 \\
Classification, extraction (2--3) & Calibration sample & After Phase
0 \\
Handwriting capture (3H) & 0.5 quality bands & After 0.5 \\
\textbf{Attribution reference table (8.2)} & ACK format spec &
\textbf{After Phase 0 --- independent of scan processing} \\
Attribution resolution (8.3) & Extraction + entity resolution & After
Phase 3 \\
Completeness (4) & Attribution, GL export & After 3A \\
Data model and load spec (6) & Phase 0 taxonomy & After Phase 0 \\
Analytics definitions (5) & Decision list & After Phase 0 \\
Email (7) & Screening agreement & Any time after Phase 6 \\
\end{longtable}
}

The attribution reference table is on the critical path and is buildable
immediately from the client's order-entry export, independent of
anything happening to the scans. Start it early; it de-risks §8 more
than any other single action.

\begin{center}\rule{0.5\linewidth}{0.5pt}\end{center}

\section{15. Outputs and Deliverables}\label{outputs-and-deliverables}

\subsection{15.1 Phase deliverables}\label{phase-deliverables}

{\def\LTcaptype{none} % do not increment counter
\begin{longtable}[]{@{}
  >{\raggedright\arraybackslash}p{(\linewidth - 2\tabcolsep) * \real{0.5000}}
  >{\raggedright\arraybackslash}p{(\linewidth - 2\tabcolsep) * \real{0.5000}}@{}}
\toprule\noalign{}
\begin{minipage}[b]{\linewidth}\raggedright
Phase
\end{minipage} & \begin{minipage}[b]{\linewidth}\raggedright
Deliverable
\end{minipage} \\
\midrule\noalign{}
\endhead
\bottomrule\noalign{}
\endlastfoot
0 & Constraint document; calibration sample; ACK/job format
specification \\
0.5 & Grayscale/bilevel split; quality distribution; JBIG2 register;
remediation recommendation \\
1 & Normalized page store; reassembled document register; deduplication
report \\
2 & Classification rules or model; accuracy report \\
3 & Extracted field dataset; consensus and arithmetic exception
queues \\
3H & Handwriting region index with cropped images; targeted
transcriptions; amendment table \\
3A & Attribution reference table; attributed event stream; coverage
report; unattributable register; selling-location assignment \\
4 & Completeness Report; reconciliation statements; matched event
ledger \\
4Q & Golden set; sampling plan; one-sided MUS overstatement bound;
separate reviewed attribute results \\
5 & Analytics layer including job and geographic dashboards \\
6 & Canonical dataset; data dictionary; ERD; load order; load
manifests \\
7 & Email-derived contacts; attribution gap-fill; attachment discrepancy
report \\
\end{longtable}
}

\subsection{15.2 Final output}\label{final-output}

A single validated corpus supporting four uses:

\begin{enumerate}
\def\labelenumi{\arabic{enumi}.}
\tightlist
\item
  \textbf{Financial analysis} --- every purchase and payment documented,
  matched, reconciled, with variances quantified.
\item
  \textbf{Job-level analysis} --- every dollar tied to an ACK or job
  number, enabling per-reference profitability across the measured
  period.
\item
  \textbf{Geographic and operational analysis} --- revenue by selling
  location; shipment, lane, carrier, and transit performance.
\item
  \textbf{CRM foundation} --- deduplicated, entity-resolved,
  load-ordered, with enforced integrity and full traceability from every
  record to a source page.
\end{enumerate}

\subsection{15.3 What the client can state on
completion}\label{what-the-client-can-state-on-completion}

\begin{itemize}
\tightlist
\item
  What the data says.
\item
  How accurate it is, with a numeric confidence bound, reported
  separately for the better and worse scan quality bands.
\item
  What is missing, enumerated.
\item
  \textbf{Which dollars could not be attributed, and why, with a dollar
  figure.}
\item
  Where every individual figure came from.
\end{itemize}

\begin{center}\rule{0.5\linewidth}{0.5pt}\end{center}

\section{16. Immediate Next Actions}\label{immediate-next-actions}

{\def\LTcaptype{none} % do not increment counter
\begin{longtable}[]{@{}
  >{\raggedright\arraybackslash}p{(\linewidth - 6\tabcolsep) * \real{0.2500}}
  >{\raggedright\arraybackslash}p{(\linewidth - 6\tabcolsep) * \real{0.2500}}
  >{\raggedright\arraybackslash}p{(\linewidth - 6\tabcolsep) * \real{0.2500}}
  >{\raggedright\arraybackslash}p{(\linewidth - 6\tabcolsep) * \real{0.2500}}@{}}
\toprule\noalign{}
\begin{minipage}[b]{\linewidth}\raggedright
\#
\end{minipage} & \begin{minipage}[b]{\linewidth}\raggedright
Action
\end{minipage} & \begin{minipage}[b]{\linewidth}\raggedright
Owner
\end{minipage} & \begin{minipage}[b]{\linewidth}\raggedright
Blocks
\end{minipage} \\
\midrule\noalign{}
\endhead
\bottomrule\noalign{}
\endlastfoot
1 & Total page count across all scan files & Client & QA sizing and cost
model \\
2 & Corpus access for the Phase 0.5 probe & Client & Quality banding,
remediation decision \\
3 & \textbf{Relationship-key format specification, including changes
over the period} & Client & Attribution and schema mapping \\
4 & \textbf{Reference-system export, if one exists} & Client &
Highest-confidence attribution and relationship validation \\
5 & \textbf{Approved reporting dimensions and business questions} &
Client & Analytics and schema design \\
6 & Confirm whether grayscale or colour masters exist for any pages &
Client & Remediation option \\
7 & Confirm whether relationship keys appear handwritten on documents &
Client & §7 scope \\
8 & Accounting system GL export & Client & Phase 4 reconciliation \\
9 & Decision list & Client & Phase 5 scope \\
10 & Email scope and screening agreement & Client & Phase 7 only --- not
blocking \\
\end{longtable}
}

Items 3, 4, and 5 materially change the shape of the work. Item 4 can
reduce the cost of attribution substantially when an authoritative
export exists; when it does not, inferred proposals remain clearly
labelled and review-bound.

\begin{center}\rule{0.5\linewidth}{0.5pt}\end{center}

\section{17. Reference Implementation}\label{reference-implementation}

A working toolkit implements the load-bearing controls in this plan. It
includes optional, client-approved OpenAI Responses, Gemini-on-Vertex
and OpenRouter page-extraction adapters, Document AI corroboration,
Google Cloud Vision handwriting proposals, and a retained
image-comparison ensemble. These are implemented integrations, not
claims of deployed accuracy or automatic approval. Additional
independent OCR/HTR providers, calibrated production detectors, broad
automatic classification, and target-system execution still depend on
client-selected integrations and acceptance. The validation and control
layer determines whether the retained output can proceed; no provider or
proposal lane may bypass it. The optional visual-intake journal and
approved-fact MCP/API surfaces are described separately in §0.3 and
§17.1.

\subsection{17.1 Scripts}\label{scripts}

This is a representative implementation map, not the complete CLI
catalogue. Use
\texttt{references/\allowbreak{}command-\allowbreak{}line-\allowbreak{}reference.md},
its parser-derived help catalogue, and the operations skill's lane table
for every current command and flag.

{\def\LTcaptype{none} % do not increment counter
\begin{longtable}[]{@{}
  >{\raggedright\arraybackslash}p{(\linewidth - 4\tabcolsep) * \real{0.3333}}
  >{\raggedright\arraybackslash}p{(\linewidth - 4\tabcolsep) * \real{0.3333}}
  >{\raggedright\arraybackslash}p{(\linewidth - 4\tabcolsep) * \real{0.3333}}@{}}
\toprule\noalign{}
\begin{minipage}[b]{\linewidth}\raggedright
Script
\end{minipage} & \begin{minipage}[b]{\linewidth}\raggedright
Implements
\end{minipage} & \begin{minipage}[b]{\linewidth}\raggedright
Plan section
\end{minipage} \\
\midrule\noalign{}
\endhead
\bottomrule\noalign{}
\endlastfoot
\texttt{scan\_\allowbreak{}profile.py} & Scan buckets, JBIG2 hazards,
quality bands, false-colour checks, and rescan evidence & §3 \\
\texttt{ingest\_\allowbreak{}pages.py} & Byte-verified source retention,
immutable one-page PDFs, relative provenance, native text, conservative
type labels, and exceptions & §4-5 \\
\texttt{reassemble\_\allowbreak{}pages.py} & Safe manifest-path
validation, review-required page groups, unordered candidates, and exact
duplicates & §4.3 \\
\texttt{preprocess\_\allowbreak{}pages.py} & Grayscale master, enhanced
grayscale, and three complementary binarized siblings & §4/6 \\
\texttt{openai\_\allowbreak{}adapter.py} & Optional structured
extraction/type/handwriting-region proposal voter with raw-response
retention & §5-7 \\
\texttt{google\_\allowbreak{}document\_\allowbreak{}ai\_\allowbreak{}adapter.py}
& Explicitly enabled independent OCR/table evidence with page hash, raw
response, token, and table-cell provenance & §4/6 \\
\texttt{table\_\allowbreak{}comprehension.py} & Source-native table
profile, evidence-linked rows, non-independent audit, exact-registry
mapping, and aggregated review/quality output & §6 \\
\texttt{table\_\allowbreak{}comprehension\_\allowbreak{}corpus.py} &
Sequential multi-manifest layout-context runner with bounded
amendment-only refinement & §6 \\
\texttt{independent\_\allowbreak{}table\_\allowbreak{}reconcile.py} &
Deterministic approved-map comparison of source rows and Document AI
cells & §6 \\
\texttt{detection\_\allowbreak{}calibration.py} & Golden-set evaluator
for handwriting-region and party-role proposals; no automatic deployment
& §5/7 \\
\texttt{llm\_\allowbreak{}adjudication.py} & Disabled-by-default,
nonfinancial printed amendment proposals after independent agreement &
§6.3 \\
\texttt{schema\_\allowbreak{}discovery.py} & Proposal-only source-label,
entity, relationship, and candidate-location discovery & §0/6 \\
\texttt{allocation\_\allowbreak{}policy.py} & Client-approved,
effective-dated allocation and sales-credit policy application &
§8/11 \\
\texttt{consensus.py} & Independent multi-engine agreement, per-field
flags, and exception routing & §6.2 \\
\texttt{arithmetic\_\allowbreak{}check.py} & Line, subtotal, charge,
total, and payment arithmetic self-proof & §6.3 \\
\texttt{adjudicate.py} & Bounded amendment proposal for a uniquely
arithmetic-supported rejected total & §6.3 \\
\texttt{validate\_\allowbreak{}extraction.py} & Deterministic
provenance, date, currency, identifier, and plausibility checks &
§6.3 \\
\texttt{address\_\allowbreak{}normalize.py} & Raw-address retention,
local CRM components, and optional capped Address Validation evidence &
§6.4 \\
\texttt{handwriting\_\allowbreak{}review.py} & Strict independent HTR
reconciliation with a two-pass cap and amendment-only financial output &
§7 \\
\texttt{entity\_\allowbreak{}resolve.py} & Party clustering with a
reversible merge log and ambiguity review & §6.6 \\
\texttt{attribution.py} & Eight-rank reference resolution, evidence
chains, location hierarchy, and unattributable register & §8 \\
\texttt{completeness.py} & Sequence/calendar gaps, GL variance, aging
closure, and explicit gate & §9 \\
\texttt{sampling.py} & Monetary-unit and stratified sampling,
projection, confidence bound, and escalation & §10 \\
\texttt{final\_\allowbreak{}review\_\allowbreak{}queue.py} & Versioned,
exhaustive eight-field canonical review package for supplied artifacts &
§3-10 \\
\texttt{operations.py} & Run manifests/state, adapter contracts, privacy
inventory, and strict reviewer exports & Cross-phase \\
\texttt{canonical\_\allowbreak{}deploy.py} & Validated non-secret
PostgreSQL schema deployment plan and explicit execution & §12 \\
\texttt{canonical\_\allowbreak{}load.py} & Ordered, checksummed load
plan with review status, batch lineage, and unique idempotency keys &
§12 \\
\texttt{csv\_\allowbreak{}api\_\allowbreak{}staging.py} & No-send CSV
files and API-operation envelope checked against the canonical load plan
& §12 \\
\texttt{retrieval\_\allowbreak{}store.py} & Approved-fact SQLite store
with document, invoice, line, job, attribution, and location context &
§12 \\
\texttt{retrieval\_\allowbreak{}mcp.py} & Approved document/CRM search,
exact lookup, schema, account cards, governed queries, sales and reports
over stdio; optional separate proposal intake & §0.3/12 \\
\texttt{retrieval\_\allowbreak{}remote\_\allowbreak{}mcp.py} & OAuth
Streamable HTTP transport, owner-bound report downloads and optional
intake JSON API & §0.3/12 \\
\texttt{retrieval\_\allowbreak{}https.py} & TLS/bearer read-only GET
business queries; not an MCP endpoint & §12 \\
\texttt{visual\_\allowbreak{}ingestion\_\allowbreak{}export.py} &
Receipt-verified source-only archive and image-PDF handoff; no proposal
promotion & §0.3/4 \\
\texttt{crm\_\allowbreak{}import\_\allowbreak{}package.py} & Verified
no-send common-object CSVs, mapping template and write plan & §12 \\
Schema-sample workbook builder & Regenerates clearly labelled fictional
schema samples & Delivery documentation \\
\end{longtable}
}

\subsection{17.2 Chain}\label{chain}

The operational Python stages use JSON contracts and carry records
forward; the review workbook and local retrieval database are derived
delivery views:

\begin{verbatim}
scan_profile   → profile.json
ingest_pages   → ingestion_manifest.json
               → classification_exceptions.json
reassemble     → reassembly.json
               → reassembly_exceptions.json
preprocess     → preprocessing_manifest.json
openai_adapter → openai_engine.json + openai_adapter.json
               → openai_exceptions.json + openai_raw/
schema_discovery → schema_discovery.json + schema_discovery_exceptions.json
allocation_policy → allocation_results.json + allocation_exceptions.json
consensus      → consensus.json  + exceptions.json
arithmetic     → proofed.json
adjudicate     → amendments.json
               → adjudication_exceptions.json
validate       → validated.json
               → validation_exceptions.json
address_normalize → address_normalized.json
               → address_exceptions.json
handwriting    → handwriting_decisions.json
               → handwriting_client_review.json
entity_resolve → parties.json    + merges.json
attribution    → attributed.json + unattributable.json
completeness   → completeness.json
sampling       → sample_plan.json
               → accuracy_statement.json
preprocess     → preprocessing_manifest.json
operations     → run_manifest.json + run_state.json
               → privacy_inventory.json
final_review   → final_client_review.json
review_export  → final_client_review.csv
               → final_client_review.html
               → final_client_review.xlsx
canonical_load → canonical_load_plan.json
retrieval_store → business_retrieval.sqlite
retrieval_mcp  → approved-fact CRM queries, reports, and lookup
\end{verbatim}

\subsection{17.3 Delivery interface and update
rule}\label{delivery-interface-and-update-rule}

The schema-sample implementation is
\texttt{scripts/\allowbreak{}build\_\allowbreak{}schema\_\allowbreak{}samples\_\allowbreak{}workbook.mjs}.

The canonical final-review package is JSON; CSV, XLSX, and HTML are
regenerated reviewer views. The XLSX workbook has a filterable
\texttt{Client\ Review} sheet plus \texttt{Header\ Definitions},
\texttt{Schema\ Definitions}, and \texttt{Runbook} sheets. Update the
canonical JSON or create an amendment, then regenerate outputs; never
use a workbook edit to alter a source value or close a gate. The
\texttt{examples/\allowbreak{}} folder contains only artefacts labelled
\textbf{SAMPLE - FICTIONAL - NOT CLIENT DATA}. It demonstrates extracted
and derived field shapes, not client evidence, calibration, or review
output.

Address components use the role-prefixed CRM names
\texttt{*\_\allowbreak{}address\_\allowbreak{}line1},
\texttt{*\_\allowbreak{}address\_\allowbreak{}line2},
\texttt{*\_\allowbreak{}city},
\texttt{*\_\allowbreak{}state\_\allowbreak{}or\_\allowbreak{}region},
\texttt{*\_\allowbreak{}postal\_\allowbreak{}code}, and
\texttt{*\_\allowbreak{}country\_\allowbreak{}code}; raw addresses
remain immutable evidence and local parsing does not prove
deliverability. When client-approved Google Address Validation is
explicitly enabled with a protected credential and a cap of at most
5,000 calls per run, its deliverability and geocode result is separately
retained evidence; uncertain, failed, exhausted, or conflicting results
remain final-review work. Never put the credential in a handoff or
delivered artifact.

Runtime/provider controls are centralized in the Git-ignored root
\texttt{.env}, copied from \texttt{.env.example}; explicit command-line
options and existing environment values take precedence.
\texttt{references/\allowbreak{}runtime-\allowbreak{}configuration.md}
defines every setting. The selected adapter records its provider/model
and supported reasoning effort, limits, timeout, retries, and non-secret
credential reference. Resolve these from the specific lane; do not infer
a model or reasoning default from another lane or from an old run.

\subsection{17.4 Verification performed}\label{verification-performed}

The chain was exercised end to end against a synthetic 60-document
corpus carrying deliberately seeded faults. Results:

{\def\LTcaptype{none} % do not increment counter
\begin{longtable}[]{@{}
  >{\raggedright\arraybackslash}p{(\linewidth - 2\tabcolsep) * \real{0.5000}}
  >{\raggedright\arraybackslash}p{(\linewidth - 2\tabcolsep) * \real{0.5000}}@{}}
\toprule\noalign{}
\begin{minipage}[b]{\linewidth}\raggedright
Seeded fault
\end{minipage} & \begin{minipage}[b]{\linewidth}\raggedright
Detected
\end{minipage} \\
\midrule\noalign{}
\endhead
\bottomrule\noalign{}
\endlastfoot
Broken printed arithmetic & Yes --- flagged as exception, handwriting
correctly ruled out \\
Handwritten quantity amendment breaking the total & Yes --- flagged,
handwriting correctly implicated \\
Two-of-three engine disagreement & Yes --- accepted with
\texttt{consensus\_\allowbreak{}2of3}, queued for review \\
Three-way engine disagreement & Yes --- hard exception, not
auto-resolved \\
Malformed ACK number (\texttt{AK-99}) & Yes --- format-rejected, then
recovered via shipment inheritance \\
Vendor name variants (\texttt{ACME\ Corp} / \texttt{ACME\ Corporation} /
\texttt{Acme\ Corp.}) & Yes --- merged to one party \\
Dotted legal suffix (\texttt{Globex\ LLC} / \texttt{Globex\ L.L.C.}) &
Yes --- merged after normalization fix \\
GL inflated 4\% against documents & Yes --- variance reported, not
plugged \\
Bilevel, grayscale, and false-colour pages & Yes --- all bucketed
correctly \\
Sampling tolerance breach & Yes --- escalation to 100\% review
triggered \\
\end{longtable}
}

The synthetic test corpus exercised four example resolution ranks
(explicit printed, explicit handwritten, cross-reference, and inherited
linkage). This is test evidence for control behavior, not a claim that
any client corpus has achieved attribution coverage.

\subsection{17.5 Production operations
controls}\label{production-operations-controls}

The reference implementation now creates immutable run manifests, hashes
non-secret configuration, blocks resumptions against a different
manifest, validates vendor-neutral OCR/HTR handoff contracts, avoids
embedded credentials, inventories obvious privacy patterns without
changing evidence, creates grayscale and an enhanced grayscale and three
complementary binarized sibling page variants, and exports the final
review package to CSV, a formatted XLSX workbook, and static HTML. CI
executes strict coverage, lint, acceptance, public-corpus validation,
and LaTeX rendering. These controls improve reproducibility and
operating safety; they do not establish an accuracy percentage without a
representative golden set.

\subsection{17.6 What remains to build}\label{what-remains-to-build}

\begin{itemize}
\tightlist
\item
  Additional independent OCR/HTR vendor adapters and calibrated local
  handwriting-region detectors
\item
  Deskew, despeckle, rotation, and scan-specific enhancement beyond the
  retained comparison ensemble
\item
  Calibrated automatic classification and party-role assignment beyond
  the high-signal rules and proposal-only model lane
\item
  Approved-template layout-delta comparison and a configured stamp
  library (§3.5)
\item
  Email ingestion, if Phase 7 proceeds (§13)
\item
  A target-specific CRM adapter, remote transport, and load
  reconciliation against the chosen system (§12)
\item
  Representative golden-set calibration before any numerical production
  accuracy claim
\end{itemize}

The control layer being complete first is deliberate. It means every
component added above is measured against working checks from the day it
lands, rather than being trusted until something downstream fails.

\begin{center}\rule{0.5\linewidth}{0.5pt}\end{center}

\section{Appendix A --- Plain-Language Illustrative
Scenario}\label{appendix-a-plain-language-illustrative-scenario}

\emph{This appendix restates the plan without technical vocabulary. It
is a non-normative example only: the actual document families, coverage
periods, relationship keys, locations, and optional evidence sources are
established in Phase 0. The implemented controls and explicit
future-work boundaries in the body of this plan govern if this appendix
and the technical sections differ.}

\subsection{A.1 What you have and what you
want}\label{a.1-what-you-have-and-what-you-want}

You have a retained collection of business records --- invoices,
shipping receipts, delivery slips, or other approved document families
--- scanned into PDF files. Some pages may have handwriting or degraded
image quality.

You want reliable, analyzable numbers; every monetary amount tied to a
configured business reference where evidence permits; clear origin or
other approved analytic dimensions; and clean party records for an
agreed target system.

You may also have a separately scoped email archive that can be
considered later as corroborating evidence.

\subsection{A.2 The core difficulty}\label{a.2-the-core-difficulty}

Run the scans through ordinary scanning software and you get a
spreadsheet full of numbers. Some are wrong. You will not know which.

That is the real danger. Wrong numbers that look right are worse than no
numbers, because you will act on them. A total off by a digit, a
customer split across three phantom accounts, an invoice scanned twice
and counted twice --- none of these announce themselves.

So the plan is built around one question: \textbf{how do we know the
numbers are right without relying on unchecked automation?}

\subsection{A.3 How we answer that}\label{a.3-how-we-answer-that}

\textbf{We compare genuinely independent readings and deterministic
controls.}

Not one system reading each page, but three independent ones, compared.
When they agree, we accept it. When they disagree, it goes on a review
list. Originally this plan called for two systems; black-and-white scans
are harder to read accurately than colour ones, so we have added a
third.

\textbf{We make every invoice prove itself with arithmetic.}

An invoice contains its own answer key. The line items should add up to
the subtotal; the subtotal plus tax and freight should equal the total.
If a document does not add up, something was read wrong --- and we know
it, even if the software was confident.

This has become the most important check in the whole project, for a
reason explained in A.5.

\textbf{We hand-check a sample, chosen to matter.}

Not a random sample. A sample weighted toward \textbf{dollars}. A
\$10,000 invoice is far more likely to be checked than a \$50 one,
because that is where the risk is. Everything above a threshold you set
gets checked for certain.

The result is not ``we spot-checked it and it looked fine.'' It is a
one-sided 95\% upper bound on monetary overstatement, paired with
separately disclosed understatement observations and completeness
results. That is a number whose scope a client can evaluate without
mistaking it for proof of missing records.

And if any category fails the check --- say invoices from one vendor in
2022 turn out error-prone --- we do not accept them with a warning
label. \textbf{We go back and check every single one.} Sampling does not
replace full review. It tells us where full review is needed.

\subsection{A.4 What black-and-white scanning costs
you}\label{a.4-what-black-and-white-scanning-costs-you}

This is worth understanding plainly, because it changes several things.

When a page is scanned in colour, blue or black pen ink is visibly a
different colour from printed black text, and finding handwriting is
easy. In black and white, pen and print become the same, and handwriting
has to be found by the \emph{shape} of the writing --- which works, but
not as reliably.

The same applies to stamps. A red ``PAID'' stamp is unmistakable in
colour. In black and white it blends into whatever it was stamped over.

What we do about it:

\begin{itemize}
\tightlist
\item
  \textbf{Read everything three times instead of twice}, as described
  above.
\item
  \textbf{Where scans are grayscale rather than pure black-and-white},
  process each page three different ways and treat each as another
  independent reading. Costs only computer time.
\item
  \textbf{Look for stamps by shape} instead of colour, using a small
  library of your recurring stamps.
\item
  \textbf{Assume handwriting might be on any page}, rather than trusting
  that we found all of it.
\item
  \textbf{Lean much harder on the arithmetic check}, which is the
  subject of A.5.
\end{itemize}

One thing worth checking: not all black-and-white scans are equal. Some
keep grey shading; some are pure black-and-white with the grey thrown
away. The second kind is meaningfully harder to work with. We will
measure which you have in the first hour. If any pages are the harder
kind and you still have the original paper or an earlier better scan,
re-scanning just those pages is the cheapest accuracy improvement
available anywhere in this project.

\subsection{A.5 About the handwriting}\label{a.5-about-the-handwriting}

You have said the handwriting can be ignored. We have accepted that,
mostly --- but there are two places where it cannot be, and it is better
to say so now than to discover it in month three.

\textbf{First, your job numbers.} You need every dollar tied to an ACK
or job number. On paperwork like this, those numbers are very often
\emph{written on by hand} --- added at order entry or at receiving. If
we throw away all handwriting and your job numbers are among it, we
cannot meet your main requirement at all. So we do need to read
handwritten job and ACK numbers, even if we read nothing else.

\textbf{Second, corrections.} When someone crosses out ``40'' and writes
``38 --- two damaged'', the printed 40 is now wrong. Ignoring the
handwriting does not make the printed number right. It makes it wrong,
and unflagged. The arithmetic check will catch most of these anyway ---
the invoice will stop adding up --- so the work does not actually go
away by ignoring it. Only our ability to resolve it quickly does.

\textbf{What we are doing instead --- and this is what your instruction
buys you:}

The target integration will \emph{find} the handwriting and retain a
crop, but not transcribe most of it. Release 1.0 already defines the
canonical region/evidence contract and accepts independent HTR results;
calibrated local detection, cropping, and CRM attachment remain
implementation work. Once connected, a user can see a handwriting flag
and its retained crop without anyone typing it out.

The future targeted HTR integration reads handwriting in four
situations: it looks like a job or ACK number; the invoice does not add
up and the handwriting is sitting on a quantity or price; it appears to
be a cheque number or payment mark; or it landed in the quality-check
sample.

That preserves nearly all of your cost saving --- the bulk of
handwriting is delivery notes and initials that genuinely do not need
transcribing --- while protecting the two things that would otherwise
break. And because the images are stored, you can change your mind later
and have them transcribed without re-processing anything.

\subsection{A.6 Tying every dollar to a
job}\label{a.6-tying-every-dollar-to-a-job}

This is a bigger job than it sounds, and it gets its own phase.

The problem: roughly half your documents will not have the job number
written on them anywhere. A carrier's freight bill has a tracking number
and a shipping reference, not your job number. So the job number has to
be \textbf{worked out} by following the trail from one document to
another.

How we work it out, in order of preference:

\begin{enumerate}
\def\labelenumi{\arabic{enumi}.}
\tightlist
\item
  It is printed on the document. Best case.
\item
  It is handwritten on the document. We read it (see A.5).
\item
  The purchase order number is on it, and we can look up which job that
  PO belongs to.
\item
  The document relates to a shipment we have already tied to a job ---
  so it inherits the job number.
\item
  It relates to an invoice we have already tied to a job --- same idea.
\item
  The customer, date, and amount uniquely match one job in your records.
\item
  An email mentions it (only if we do the email phase).
\item
  Someone assigns it by hand, and we record who and when.
\item
  \textbf{It genuinely cannot be attributed} --- and we record that
  explicitly, with a reason and a dollar amount.
\end{enumerate}

Steps 4 and 5 are what make this achievable. They are how a freight bill
that mentions no job number still ends up correctly filed.

\textbf{The most useful thing you can give us:} if your order-entry or
quoting system can export a list of every ACK number with its date,
customer, and amount, that list becomes the master reference and makes
this entire phase far cheaper and far more reliable. This is the single
highest-value item on the list of things we need from you.

\textbf{What ``done'' looks like:} not ``97\% of dollars attributed.''
It is that 100\% of dollars are either attributed \emph{or} explicitly
listed as unattributable with a reason and an amount. Ninety-seven
percent with the other three percent unexplained is not a passing result
--- the missing three percent is exactly what someone will ask about.

\subsection{A.7 Which city the sale came
from}\label{a.7-which-city-the-sale-came-from}

You want to know what city each sale came from. Before we build
anything, we need to pin down what that means, because it has four
possible meanings and they give different answers:

\begin{itemize}
\tightlist
\item
  \textbf{Which of your offices sold it} --- a branch or territory
  performance question
\item
  \textbf{Where the goods shipped from} --- a warehouse question
\item
  \textbf{Where the customer is} --- a market question
\item
  \textbf{Where the goods went} --- a destination question
\end{itemize}

Most people asking this question mean the first one. It is also the
hardest to find, because it is often not written on the document at all
--- it has to be worked out from the letterhead, from the ACK number
prefix if your numbering encodes location, or from which salesperson
handled it.

We will capture all four where the data supports it, so if you decide
later that you meant a different one, nothing has to be redone. But tell
us which one you mean now, because it shapes what we build.

And every city figure will show how it was worked out. A city total that
is 60\% guessed from ACK prefixes is a different quality of number than
one that is written on the documents, and you should be able to see
which you are looking at.

\subsection{A.8 Making sure nothing is
missing}\label{a.8-making-sure-nothing-is-missing}

Accuracy is whether the numbers we captured are right.
\textbf{Completeness} is whether we captured all of them --- harder,
because a document that was never scanned leaves no trace.

\textbf{Numbering gaps.} Invoices run in sequence. If we have 1001,
1002, and 1004, then 1003 exists somewhere and we do not have it. We
list every gap. We now also do this on your ACK numbers --- and since
those are your own numbers, gaps in them are much stronger evidence than
gaps in a supplier's.

\textbf{Calendar gaps.} We count documents by supplier by month. If a
supplier you use every month has nothing in March 2023, a batch of scans
is probably missing.

\textbf{Comparison to your books.} We compare our totals against your
accounting system, period by period. If your books say \$2.4 million in
2022 and our documents account for \$2.1 million, there is \$300,000 of
paperwork we do not have. \textbf{We report that. We never quietly
adjust the numbers to make them match.}

\textbf{Loose ends.} Every invoice should end somewhere --- paid,
credited, written off, or still open. Anything that just stops gets
listed.

\textbf{Unattributed dollars.} Per A.6.

All of this goes into a report that travels with the data. Every chart
you get shows what portion of the records it is based on. If a chart is
built on 94\% of documents, the chart says so.

\subsection{A.9 Cleaning up customer
names}\label{a.9-cleaning-up-customer-names}

Small-sounding, and it determines whether the CRM works.

Over time the same customer may appear as ``ACME Corp'', ``Acme
Corporation'', ``ACME CORP.'', and probably ``Acme Corp - Chicago''.
These are one customer. If they arrive in your new CRM as four accounts,
then every question you want to ask --- biggest customers, who stopped
buying, which accounts are growing --- gives a wrong answer. You would
see four mid-sized accounts instead of one large one.

We merge these into one clean record per real business, keeping a log of
every decision so it can be reviewed and undone.

\subsection{A.10 The email option}\label{a.10-the-email-option}

An email archive is an optional, separately authorized step. Deferring
it costs nothing, and the design leaves room so nothing has to be
rebuilt if it is later approved.

\textbf{One caveat first:} email may cover only part of the document
period. Any number that draws on it must state the overlap window ---
otherwise a chart can show what looks like a business change when it is
only an evidence-coverage boundary.

\textbf{Why it is worth considering, in order:}

\textbf{Attachments.} This is the big one. Emails often contain the
original, clean, computer-generated PDF of an invoice that we only have
as a degraded black-and-white scan. Where we find one, the attachment is
\emph{better evidence than the scan} and use it as a separately retained
source. For a degraded scan corpus, this alone may justify the phase.
Comparing the clean version to the scanned version measures the
scan-quality effect without overwriting either source.

\textbf{Filling in missing job numbers.} For dollars we could not
attribute from documents alone, an email is often the only remaining
evidence.

\textbf{Missing payment records.} Remittance notices and payment
confirmations close out invoices that otherwise look unpaid.

\textbf{Contact details for the CRM.} Names, titles, phone numbers,
current email addresses --- things your paper documents barely contain
and your CRM genuinely needs.

\textbf{The story behind the numbers.} Quote, change order, dispute,
settlement --- the documents record the outcome; the emails explain it.

\textbf{One firm rule:} an email never overrides a document. If an email
says an invoice was for \$5,500 and the invoice says \$5,507.50, that is
a discrepancy for a person to resolve, not a correction to apply. The
exception is an actual PDF attachment, which is a document, not an
email.

\textbf{Before we touch any of it}, we agree in writing which mailboxes,
what date range, and what gets filtered out --- personal correspondence,
HR matters, anything privileged. That agreement comes first, not after.

\subsection{A.11 What you get at the
end}\label{a.11-what-you-get-at-the-end}

\textbf{Reliable numbers}, with a precise monetary statement --- the
one-sided 95\% upper bound on overstatement is X dollars. Observed
understatement, missing documents, and open exceptions are reported
separately, along with results by scan-quality stratum so you can see
what the scan quality cost.

\textbf{A stated list of what is missing} --- numbering gaps, calendar
gaps, and the difference between our totals and your accounting system.

\textbf{Every dollar tied to a job or ACK number}, or explicitly listed
as unattributable with a reason and an amount.

\textbf{Job-level analysis} --- what each job actually made, what it
cost, how long it took from order to payment, and which ones lost money.

\textbf{Revenue by city}, showing how each figure was determined.

\textbf{Business analysis}, including: which customers generate the most
revenue and how concentrated you are; which customers have quietly
stopped ordering; what you spend by carrier and route; where surcharges
and extra fees are accumulating; which carriers deliver on time;
duplicate invoices you may have paid twice; whether you are being
charged your contracted rates; and how your business moves seasonally.

\textbf{Clean data for your CRM}, in the correct loading order, with the
rules your developers need to keep it clean.

\textbf{Handwriting captured as images}, attached to the right
documents, flagged and searchable --- even though most of it was never
transcribed.

\textbf{Complete traceability.} Every number, anywhere in the system,
traces back to a specific page of a specific scanned document. If anyone
questions a figure, you can produce the piece of paper it came from.

\subsection{A.12 What we need from you to
start}\label{a.12-what-we-need-from-you-to-start}

\begin{enumerate}
\def\labelenumi{\arabic{enumi}.}
\tightlist
\item
  \textbf{How many pages in total?} Determines the size and cost of the
  quality checking.
\item
  \textbf{What do your ACK and job numbers look like?} Format, length,
  prefixes, and whether the convention changed over the measured period.
\item
  \textbf{Can your order-entry system export a list of ACK numbers with
  dates, customers, and amounts?} The most valuable single thing on this
  list.
\item
  \textbf{Which meaning of ``city the sale came from'' do you want?} See
  A.7.
\item
  \textbf{Do your job numbers get handwritten onto documents?}
  Determines how much handwriting work is unavoidable.
\item
  \textbf{Access to the scan files}, for the one-hour quality check.
\item
  \textbf{Do you still have original paper, or better-quality earlier
  scans?} Determines whether targeted re-scanning is an option.
\item
  \textbf{An export from your accounting system}, to reconcile against.
\item
  \textbf{A list of the questions you want answered}, in plain
  sentences. Defines what ``finished'' means.
\item
  \textbf{Email scope}, only when and if you want the email phase.
\end{enumerate}

\subsection{A.13 What has already been
built}\label{a.13-what-has-already-been-built}

The checking machinery described above is not a proposal --- it exists
and has been tested. The control toolkit implements intake, comparison
between independent readings, arithmetic proof, customer-name
resolution, job-number attribution, completeness checks, sampling,
client review packaging, and a vendor-neutral post-review load and
retrieval boundary.

They were tested against a synthetic set of sixty invoices with
deliberate faults planted in them --- a broken total, a handwritten
correction, disagreeing readings, a malformed job number, duplicate
customer names spelled four ways, and books inflated by 4\% against the
paperwork. Every planted fault was caught.

What remains client-specific is integration and authorization:
additional independent reading services, calibrated handwriting and
party-role detectors, scan-specific cleanup, a target CRM adapter, and a
representative client-approved golden set. The repository now includes
the evaluator for that golden set, exact template drift controls, and a
proposal-only decision compiler with operator authorization; those
controls do not supply the missing client evidence or apply production
facts. The current reassembly and vision lanes remain review proposals.
That order means every added component is measured against working
controls before production use.

\subsection{A.14 The one-paragraph
version}\label{a.14-the-one-paragraph-version}

We will preserve and compare independent readings of scanned documents,
make every invoice prove itself by checking its own arithmetic,
photograph and flag every piece of handwriting while only transcribing
the parts that carry job numbers or affect the money, work out a job or
ACK number for every single dollar and explicitly list the ones we
cannot, determine which of your offices each sale came from and show how
we determined it, hand-check a dollar-weighted sample and fully review
any category that fails, tell you plainly what is missing rather than
hiding it, merge duplicate customer names into single clean records, and
hand you a dataset where every number traces back to a specific page ---
with an honest, numeric statement of how accurate and how complete it
is. Your email can be folded in later to fill gaps and, where it
contains clean original PDFs of documents we only have as poor scans, to
improve on the scans themselves.

\end{document}
